\section{AIME I 2019}
        \begin{enumerate}
        	\item (\textit{AIME I 2019}) Consider the integer 
 \[N = 9 + 99 + 999 + 9999 + \cdots + \underbrace{99\ldots 99}_\text{321 digits}.\]Find the sum of the digits of $N$.


 
 \textbf{Answer: }342



\textbf{Solution 1}

Let's express the number in terms of $10^n$. We can obtain $(10-1)+(10^2-1)+(10^3-1)+\cdots+(10^{321}-1)$. By the commutative and associative property, we can group it into $(10+10^2+10^3+\cdots+10^{321})-321$. We know the former will yield $1111....10$, so we only have to figure out what the last few digits are. There are currently $321$ 1's. We know the last four digits are $1110$, and that the others will not be affected if we subtract $321$. If we do so, we get that $1110-321=789$. This method will remove three $1$'s, and add a $7$, $8$ and $9$. Therefore, the sum of the digits is $(321-3)+7+8+9=\boxed{342}$.

 -eric2020-another Eric in 2020

 


 A similar and simpler way to consider the initial manipulations is to observe that adding $1$ to each term results in $(10+100+... 10^{320}+10^{321})$. There are $321$ terms, so it becomes $11...0$, where there are $322$ digits in $11...0$. Then, subtract the $321$ you initially added. 

 





\textbf{Solution 2}

We can see that $9=9$, $9+99=108$, $9+99+999=1107$, all the way to ten nines when we have $11111111100$. Then, when we add more nines, we repeat the same process, and quickly get that the sum of digits is $\boxed{342}$ since we have to add $9\lfloor \log 321 \rfloor$ to the sum of digits, which is $9\lceil \frac{321}9 \rceil$.

 





\textbf{Solution 3 (Pattern)}

Observe how adding results in the last term but with a $1$ concatenated in front and also a $1$ subtracted ($09$, $108$, $1107$, $11106$). Then for any index of terms, $n$, the sum is $11...10-n$, where the first term is of length $n+1$. Here, that is $\boxed{342}$.

 





\textbf{Solution 4 (Official MAA)}

Write 
 \begin{align*}N &=(10-1)+(10^2-1)+\cdots+(10^{321}-1)\\ &=10+10^2+10^3+10^4+10^5+10^6+\cdots 10^{321}-321 \\ &=1110-321+10^4+10^5+10^6+\cdots+10^{321}\\ &=789+10^4+10^5+10^6+\cdots+10^{321}\\ \end{align*}
The sum of the digits of $N$ is therefore equal to $7+8+9+(321-3)=342$.

 


 




	\item (\textit{AIME I 2019}) Jenn randomly chooses a number $J$ from $1, 2, 3,\ldots, 19, 20$. Bela then randomly chooses a number $B$ from $1, 2, 3,\ldots, 19, 20$ distinct from $J$. The value of $B - J$ is at least $2$ with a probability that can be expressed in the form $\frac{m}{n}$, where $m$ and $n$ are relatively prime positive integers. Find $m+n$.


 
 \textbf{Answer: }029



\textbf{Solution 1}

$B-J \ne 0$ because $B \ne J$, so the probability that $B-J < 0$ is $\frac{1}{2}$ by symmetry.

 The probability that $B-J = 1$ is $\frac{19}{20 \times 19} = \frac{1}{20}$ because there are 19 pairs: $(B,J) = (2,1), \ldots, (20,19)$.

 The probability that $B-J \ge 2$ is $1-\frac{1}{2}-\frac{1}{20} = \frac{9}{20} \implies \boxed{029}$

 





\textbf{Solution 2}

By symmetry, the desired probability is equal to the probability that $J - B$ is at most $-2$, which is $\frac{1-P}{2}$ where $P$ is the probability that $B$ and $J$ differ by $1$ (no zero, because the two numbers are distinct). There are $20 \cdot 19 = 380$ total possible combinations of $B$ and $J$, and $1 + 18 \cdot 2 + 1 = 38$ ones that form $P$, so $P = \frac{38}{380} = \frac{1}{10}$. Therefore the answer is $\frac{9}{20} \rightarrow \boxed{029}$.

 





\textbf{Solution 3}

This problem is essentially asking how many ways there are to choose $2$ distinct elements from a $20$ element set such that no $2$ elements are adjacent. Using the well-known formula $\dbinom{n-k+1}{k}$, there are $\dbinom{20-2+1}{2} = \dbinom{19}{2} = 171$ ways. Dividing $171$ by $380$, our desired probability is $\frac{171}{380} = \frac{9}{20}$. Thus, our answer is $9+20=\boxed{029}$.-Fidgetboss\_4000

 





\textbf{Solution 4}

Create a grid using graph paper, with $20$ columns for the values of $J$ from $1$ to $20$ and $20$ rows for the values of $B$ from $1$ to $20$. Since $B$ cannot equal $J$, we cross out the diagonal line from the first column of the first row to the twentieth column of the last row. Now, since $B - J$ must be at least $2$, we can mark the line where $B - J = 2$. Now we sum the number of squares that are on this line and below it. We get $171$. Then we find the number of total squares, which is $400 - 20 = 380$. Finally, we take the ratio $\frac{171}{380}$, which simplifies to $\frac{9}{20}$. Our answer is $9+20=\boxed{029}$.

 





\textbf{Solution 5}

We can see that if $B$ chooses $20$, $J$ can choose from $1$ through $18$ such that $B-J\geq 2$. If $B$ chooses $19$, $J$ has choices $1$~$17$. By continuing this pattern, $B$ will choose $3$ and $J$ will have $1$ option. Summing up the total, we get $18+17+\cdots+1$ as the total number of solutions. The total amount of choices is $20\times19$ (B and J must choose different numbers), so the probability is $\frac{18\cdot19\div2}{20\cdot19}=\frac{9}{20}$. Therefore, the answer is $9+20=\boxed{029}$

 -eric2020

 





\textbf{Solution 6}

Similar to solution 4, we can go through the possible values of $J$ to find all the values of $B$ that makes $B-J\geq 2$. If $J$ chooses $1$, then $B$ can choose anything from $3$ to $20$. If $J$ chooses $2$, then $B$ can choose anything from $4$ to $20$. By continuing this pattern, we can see that there is $18+17+\cdots+1$ possible solutions. The amount of solutions is, therefore, $\frac{18\cdot19}{2}=171$. Now, because $B$ and $J$ must be different, we have $20\times19=380$ possible choices, so the probability is $\frac{171}{380}=\frac{9}{20}$. Therefore, the final answer is $9+20=\boxed{029}$

 -josephwidjaja

 





\textbf{Solution 7 (Official MAA)}

There are $\tbinom{20}{2}=190$ equally likely pairs $\{J,B\}$. In $19$ of these pairs $(\{1,2\},\{2,3\},\{3,4\}\dots,\{19,20\})$, the numbers differ by less than 2, so the probability that the numbers differ by at least 2 is $1-\frac{19}{190}=\frac9{10}$. Then $B-J\ge 2$ holds in exactly half of these cases, so it has probability $\frac12\cdot\frac9{10}=\frac{9}{20}$. The requested sum is $9+20=\boxed{029}$.

 





\textbf{Solution 8 (Compliment Count)}

Impose the compliment.

 There are $20 \cdot 19 = 380$ possible ordered pairs $(J,B)$. We wish to find all $(J,B):~B - J < 2$. That only happens if either $B - J = 1$ or $B < J$, but never $B = J$ as that contradicts $B \neq J \in \{1,2,3,\dots,20\}$. Then, for $J=1$, $B = 2$, for $J = 2$, $B = 1,3$, and for all $J = k~:~1\le k \le 19$, if $C(B)$ denotes the number of possible $B$ given $J$, $C(B) = k$. Then, for the special $k=20$, we see only 19 values works.

 The requested compliment is then $\frac{1+2+3+\dots+19 +19}{380} = \frac{209}{380} = \frac{11}{20}$. The requested probability is $1 - \frac{11}{20} = \frac{9}{20} \implies \boxed{029}$.

 




	\item (\textit{AIME I 2019}) In $\triangle PQR$, $PR=15$, $QR=20$, and $PQ=25$. Points $A$ and $B$ lie on $\overline{PQ}$, points $C$ and $D$ lie on $\overline{QR}$, and points $E$ and $F$ lie on $\overline{PR}$, with $PA=QB=QC=RD=RE=PF=5$. Find the area of hexagon $ABCDEF$.


 
 \textbf{Answer: }120



\textbf{Solution 1}

We know the area of the hexagon $ABCDEF$ to be $\triangle PQR- \triangle PAF- \triangle BCQ- \triangle RED$. Since $PR^2+RQ^2=PQ^2$, we know that $\triangle PRQ$ is a right triangle. Thus the area of $\triangle PQR$ is $150$. Since $\triangle PQR$ is a right triangle, we also know that $\sin \angle PQR = \frac{15}{25} = \frac{3}{5}$. Then the area of $\triangle BCQ = \frac12 \cdot BQ \cdot CQ \cdot \sin \angle PQR= \frac{25}{2}\cdot \frac{3}{5}=\frac{15}{2}$. Preceding in a similar fashion for $\triangle PAF$, the area of $\triangle PAF$ is $10$. Since $\angle ERD = 90^{\circ}$, the area of $\triangle RED=\frac{25}{2}$. Thus our desired answer is $150-\frac{15}{2}-10-\frac{25}{2}=\boxed{120}$

 





\textbf{Solution 2}

Let $R$ be the origin. Noticing that the triangle is a 3-4-5 right triangle, we can see that $A=(4,12), B=(16,3), C=(15,0), D=(5,0), E=(0,5)$, and $F=(0,10)$. Using the shoelace theorem, the area is $\boxed{120}$.Shoelace theorem:Suppose the polygon $P$ has vertices $(a_1, b_1)$, $(a_2, b_2)$, ... , $(a_n, b_n)$, listed in clockwise order. Then the area of $P$ is

 
 \[\dfrac{1}{2} |(a_1b_2 + a_2b_3 + \cdots + a_nb_1) - (b_1a_2 + b_2a_3 + \cdots + b_na_1)|\]You can also go counterclockwise order, as long as you find the absolute value of the answer.

 .

 





\textbf{Solution 3 (Easiest, uses only basic geometry too)}

Note that $\triangle{PQR}$ has area $150$ and is a $3$-$4$-$5$ right triangle. Then, by similar triangles, the altitude from $B$ to $QC$ has length $3$ and the altitude from $A$ to $FP$ has length $4$, so $[QBC]+[DRE]+[AFP]=\frac{15}{2}+\frac{25}{2}+\frac{20}{2}=30$, meaning that $[ABCDEF]=\boxed{120}$. -Stormersyle

 





\textbf{Solution 4}

Knowing that $\triangle{PQR}$ has area $150$ and is a $3$-$4$-$5$ triangle, we can find the area of the smaller triangles $\triangle{DRE}$, $\triangle{APF}$, and $\triangle{CQB}$ and subtract them from $\triangle{PQR}$ to obtain our answer. First off, we know $\triangle{DRE}$ has area $12.5$ since it is a right triangle. To the find the areas of $\triangle{APF}$ and $\triangle{CQB}$ , we can use Law of Cosines ($c^2 = a^2 + b^2 - 2ab\cos C$) to find the lengths of $AF$ and $CB$, respectively. Computing gives $AF = \sqrt{20}$ and $CB = \sqrt{10}$. Now, using Heron's Formula, we find $\triangle{APF} = 10$ and $\triangle{CQB} = 7.5$. Adding these and subtracting from $\triangle{PQR}$, we get $150 - (10 + 7.5 + 12.5) = \boxed{120}$ -Starsher

 





\textbf{Solution 5 (Official MAA)}

Triangle $PQR$ is a right triangle with are $\frac12\cdot15\cdot20=150$. Each of $\triangle PAF,\triangle QCB,$ and $\triangle RED$ shares an angle with $\triangle PQR$. Because the area of a triangle with sides $a,\,b,$ and included angle $\gamma$ is $\frac12a\cdot b\cdot \sin\gamma,$ it follows that the areas of $\triangle PAF,\triangle QCB,$ and $\triangle RED$ are each $5\cdot5\cdot\frac{150}{ab},$ where $a$ and $b$ are the lengths of the sides of $\triangle PQR$ adjacent to the shared angle. Thus the sum of the areas of $\triangle PAF,\triangle QCB,$ and $\triangle RED$ is 
 \[5\cdot5\cdot\frac{150}{15\cdot25}+5\cdot5\cdot\frac{150}{25\cdot20}+5\cdot5\cdot\frac{150}{20\cdot15}=25\left(\frac25+\frac3{10}+\frac12\right)=30.\] Therefore $ABCDEF$ has area $150-30=\boxed{120}$.

 





\textbf{Solution 6 (Using simple trigonometry)}

Let's say that $\angle APF$ = $u$. Then, we know that $\sin(u) = 20/25 = 4/5$. Therefore, the area of $\triangle{APF}$ is $0.5*5*5*\sin(u)=10$. Now, let's do the same thing with $\triangle{BQC}$. If we name $\angle BQC$ as $w$, we know that $\sin(w) = 15/25 = 3/5$. Therefore, the area of $\triangle{BQC}$ is $0.5*5*5*\sin(w)=7.5$. Since $\triangle{ERD}$ is a right-angled isosceles triangle, the area of $\triangle{ERD} = 12.5$. In conclusion, the area of $ABCDEF$ is $\triangle{PRQ}-(\triangle{APF}+\triangle{BQC}+\triangle{ERD})=150-(10+7.5+12.5)=\boxed{120}$.

 




	\item (\textit{AIME I 2019}) A soccer team has $22$ available players. A fixed set of $11$ players starts the game, while the other $11$ are available as substitutes. During the game, the coach may make as many as $3$ substitutions, where any one of the $11$ players in the game is replaced by one of the substitutes. No player removed from the game may reenter the game, although a substitute entering the game may be replaced later. No two substitutions can happen at the same time. The players involved and the order of the substitutions matter. Let $n$ be the number of ways the coach can make substitutions during the game (including the possibility of making no substitutions). Find the remainder when $n$ is divided by $1000$.


 
 \textbf{Answer: }122



\textbf{Solution 1 (Recursion)}

There are $0-3$ substitutions. The number of ways to sub any number of times must be multiplied by the previous number. This is defined recursively. The case for $0$ subs is $1$, and the ways to reorganize after $n$ subs is the product of the number of new subs ($12-n$) and the players that can be ejected ($11$). The formula for $n$ subs is then $a_n=11(12-n)a_{n-1}$ with $a_0=1$. 

 Summing from $0$ to $3$ gives $1+11^2+11^{3}\cdot 10+11^{4}\cdot 10\cdot 9$. Notice that $10+9\cdot11\cdot10=10+990=1000$. Then, rearrange it into $1+11^2+11^3\cdot (10+11\cdot10\cdot9)= 1+11^2+11^3\cdot (1000)$. When taking modulo $1000$, the last term goes away. What is left is $1+11^2=\boxed{122}$.

 





\textbf{Solution 2 (Casework)}

We can perform casework. Call the substitution area the "bench".

 $\textbf{Case 1}$: No substitutions. There is $1$ way of doing this: leaving everybody on the field.

 $\textbf{Case 2}$: One substitution. Choose one player on the field to sub out, and one player on the bench. This corresponds to $11\cdot 11 = 121$.

 $\textbf{Case 3}$: Two substitutions. Choose one player on the field to sub out, and one player on the bench. Once again, this is $11\cdot 11$. Now choose one more player on the field to sub out and one player on the bench that was not the original player subbed out. This gives us a total of $11\cdot 11\cdot 11\cdot 10 = 13310\equiv 310 \bmod{1000}$.

 $\textbf{Case 4}$: Three substitutions. Using similar logic as $\textbf{Case 3}$, we get $(11\cdot 11)\cdot (11\cdot 10)\cdot (11\cdot 9)$. The resulting number ends in $690$.

 Therefore, the answer is $1 + 121 + 310 + 690 \equiv \boxed{122} \pmod{1000}$.

 





\textbf{Solution 3 (Official MAA)}

There is $1$ way of making no substitutions to the starting lineup. If the coach makes exactly one substitution, this can be done in $11^2$ ways. Two substitutions can happen in $11^2\cdot11\cdot 10$ ways. Similarly, three substitutions can happen $11^2\cdot11\cdot10\cdot11\cdot9$ ways. The total number of possibilities is $1+11^2+11^2\cdot11\cdot10+11^2\cdot11\cdot10\cdot11\cdot9=122+11^3(10+990)\equiv 122\pmod{1000}.$

 





\textbf{Solution 4 (No time left)}

If you make no substitutions, there is $1$ way of doing that. If you make $1$ substitution, you choose $1$ substitute and $1$ player, for a total of $11\times11=121$ ways. Then you suddenly realized you got no time, and you decided to put what you got so far, which is $1+121=\boxed{122}$, and it turned out to be the correct answer.

 




	\item (\textit{AIME I 2019}) A moving particle starts at the point $(4,4)$ and moves until it hits one of the coordinate axes for the first time. When the particle is at the point $(a,b)$, it moves at random to one of the points $(a-1,b)$, $(a,b-1)$, or $(a-1,b-1)$, each with probability $\frac{1}{3}$, independently of its previous moves. The probability that it will hit the coordinate axes at $(0,0)$ is $\frac{m}{3^n}$, where $m$ and $n$ are positive integers, and $m$ is not divisible by $3$. Find $m + n$.


 
 \textbf{Answer: }252



\textbf{Solution 1}

One could recursively compute the probabilities of reaching $(0,0)$ as the first axes point from any point $(x,y)$ as 
 \[P(x,y) = \frac{1}{3} P(x-1,y) + \frac{1}{3} P(x,y-1) + \frac{1}{3} P(x-1,y-1)\] for $x,y \geq 1,$ and the base cases are$P(0,0) = 1, P(x,0) = P(y,0) = 0$ for any $x,y$ not equal to zero.We then recursively find $P(4,4) = \frac{245}{2187}$ so the answer is $245 + 7 = \boxed{252}$.

 


 





\textbf{Solution 2}

Obviously, the only way to reach (0,0) is to get to (1,1) and then have a $\frac{1}{3}$ chance to get to (0,0). Let x denote a move left 1 unit, y denote a move down 1 unit, and z denote a move left and down one unit each. The possible cases for these moves are $(x,y,z)=(0,0,3),(1,1,2),(2,2,1)$ and $(3,3,0)$. This gives a probability of $1 \cdot \frac{1}{27} + \frac{4!}{2!} \cdot \frac{1}{81} + \frac{5!}{2! \cdot 2!} \cdot \frac{1}{243} +\frac{6!}{3! \cdot 3!} \cdot \frac{1}{729}=\frac{245}{729}$ to get to $(1,1)$. The probability of reaching $(0,0)$ is $\frac{245}{3^7}$. This gives $245+7=\boxed{252}$. 

 





\textbf{Solution 3}

Since the particle stops at one of the axes, we know that the particle must pass through $(1,1)$. Thus, it suffices to consider the probability our particle will reach $(1,1)$. Then the only ways to get to $(1,1)$ from $(4,4)$ are the following:

 (1) 3 moves diagonally    

 (2) 2 moves diagonally, 1 move left, 1 move down 

 (3) 1 move diagonally, 2 moves left and 2 moves down. 

 (4) 3 moves left, 3 moves down.

 The probability of (1) is $\frac{1}{3^3}$. The probability of (2) is $\frac{\frac{4!}{2!}}{3^4} = \frac{12}{3^4}$. The probability of (3) is $\frac{\frac{5!}{2!2!}}{3^5} = \frac{30}{3^5}$. The probability of (4) is $\frac{\frac{6!}{3!3!}}{3^6} = \frac{20}{3^6}$. Adding all of these together, we obtain a total probability of $\frac{245}{3^6}$ that our particle will hit $(1,1)$. Trivially, there is a $\frac{1}{3}$ chance our particle will hit $(0,0)$ from $(1,1)$. So our final probability will be $\frac{245}{3^6} \cdot \frac{1}{3} = \frac{245}{3^7} \implies m = 245, n = 7 \implies \boxed{252}$ 

 





\textbf{Solution 4 (Official MAA)}

All paths that first hit the axes at the origin must pass through the point $(1,1)$. There are $63$ paths from the point $(4,4)$ to the point $(1,1)$: $\tbinom{6}{3}=20$ that take $3$ steps left and $3$ steps down, $\tbinom{5}{2\,2\,1}=30$ that take $2$ steps left, $2$ steps down, and $1$ diagonal step, $\tbinom{4}{1\,1\,2}=12$ steps that take $1$ step left, $1$ steps down, and $2$ diagonal steps, and $1$ that takes $3$ diagonal steps. The total probability of moving from $(4,4)$ to $(1,1)$ is therefore 
 \[\frac{1}{3^6}\cdot20+\frac{1}{3^5}\cdot30+\frac{1}{3^4}\cdot12+\frac{1}{3^3}\cdot1=\frac{245}{3^6}.\] Multiplying by $\frac13$ gives $\frac{245}{3^7},$ the probability that the path first reaches the axes at the origin. The requested sum is $245+7=252$.

 




	\item (\textit{AIME I 2019}) In convex quadrilateral $KLMN$, side $\overline{MN}$ is perpendicular to diagonal $\overline{KM}$, side $\overline{KL}$ is perpendicular to diagonal $\overline{LN}$, $MN = 65$, and $KL = 28$. The line through $L$ perpendicular to side $\overline{KN}$ intersects diagonal $\overline{KM}$ at $O$ with $KO = 8$. Find $MO$.


 
 \textbf{Answer: }090



\textbf{Solution 1 (Trig)}

Let $\angle MKN=\alpha$ and $\angle LNK=\beta$. Let $P$ be the project of $L$ onto line $NK$. Note $\angle KLP=\beta$. 

 Then, $KP=28\sin\beta=8\cos\alpha$.Furthermore, $KN=\frac{65}{\sin\alpha}=\frac{28}{\sin\beta} \Rightarrow 65\sin\beta=28\sin\alpha$.

 Dividing the equations gives
 \[\frac{65}{28}=\frac{28\sin\alpha}{8\cos\alpha}=\frac{7}{2}\tan\alpha\Rightarrow \tan\alpha=\frac{65}{98}\]

 Thus, $MK=\frac{MN}{\tan\alpha}=98$, so $MO=MK-KO=\boxed{090}$.

 





\textbf{Solution 2 (Similar triangles)}


                    \begin{figure}[H]
                        \centering
                        \includegraphics[width=0.5\textwidth]{Solution/AIME_2019_I_6_Solution_Figure_1.png}
                    \end{figure}
                    

 First, let $P$ be the intersection of $LO$ and $KN$ as shown above. Note that $m\angle KPL = 90^{\circ}$ as given in the problem. Since $\angle KPL \cong \angle KLN$ and $\angle PKL \cong \angle LKN$, $\triangle PKL \sim \triangle LKN$ by AA similarity. Similarly, $\triangle KMN \sim \triangle KPO$. Using these similarities we see that
 \[\frac{KP}{KL} = \frac{KL}{KN}\]
 \[KP = \frac{KL^2}{KN} = \frac{28^2}{KN} = \frac{784}{KN}\]and
 \[\frac{KP}{KO} = \frac{KM}{KN}\]
 \[KP = \frac{KO \cdot KM}{KN} = \frac{8\cdot KM}{KN}\]Combining the two equations, we get
 \[\frac{8\cdot KM}{KN} = \frac{784}{KN}\]
 \[8 \cdot KM = 28^2\]
 \[KM = 98\]Since $KM = KO + MO$, we get $MO = 98 -8 = \boxed{090}$.

 Solution by vedadehhc

 





\textbf{Solution 3 (Similar triangles, orthocenters)}

Extend $KL$ and $NM$ past $L$ and $M$ respectively to meet at $P$. Let $H$ be the intersection of diagonals $KM$ and $LN$ (this is the orthocenter of $\triangle KNP$).

 As $\triangle KOL \sim \triangle KHP$ (as $LO \parallel PH$, using the fact that $H$ is the orthocenter), we may let $OH = 8k$ and $LP = 28k$.

 Then using similarity with triangles $\triangle KLH$ and $\triangle KMP$ we have

 
 \[\frac{28}{8+8k} = \frac{8+8k+HM}{28+28k}\]

 Cross-multiplying and dividing by $4+4k$ gives $2(8+8k+HM) = 28 \cdot 7 = 196$ so $MO = 8k + HM = \frac{196}{2} - 8 = \boxed{090}$. (Solution by scrabbler94)

 





\textbf{Solution 4 (Algebraic Bashing)}

First, let $P$ be the intersection of $LO$ and $KN$. We can use the right triangles in the problem to create equations. Let $a=NP, b=PK, c=NO, d=OM, e=OP, f=PC,$ and $g=NC.$ We are trying to find $d.$ We can find $7$ equations. They are
 \[4225+d^2=c^2,\]
 \[4225+d^2+16d+64=a^2+2ab+b^2,\]
 \[a^2+e^2=c^2,\]
 \[b^2+e^2=64,\]
 \[b^2+e^2+2ef+f^2=784,\]
 \[a^2+e^2+2ef+f^2=g^2,\]and 
 \[g^2+784=a^2+2ab+b^2.\]We can subtract the fifth equation from the sixth equation to get $a^2-b^2=g^2-784.$ We can subtract the fourth equation from the third equation to get $a^2-b^2=c^2-64.$ Combining these equations gives $c^2-64=g^2-784$ so $g^2=c^2+720.$ Substituting this into the seventh equation gives $c^2+1504=a^2+2ab+b^2.$ Substituting this into the second equation gives $4225+d^2+16d+64=c^2+1504$. Subtracting the first equation from this gives $16d+64=1504.$ Solving this equation, we find that $d=\boxed{090}.$(Solution by DottedCaculator)

 





\textbf{Solution 5 (5-second PoP)}


                    \begin{figure}[H]
                        \centering
                        \includegraphics[width=0.5\textwidth]{Solution/AIME_2019_I_6_Solution_Figure_2.png}
                    \end{figure}
                    Notice that $KLMN$ is inscribed in the circle with diameter $\overline{KN}$ and $XOMN$ is inscribed in the circle with diameter $\overline{ON}$. Furthermore, $(XLN)$ is tangent to $\overline{KL}$. Then, 
 \[KO\cdot KM=KX\cdot KN=KL^2\implies KM=\frac{28^2}{8}=98,\]and $MO=KM-KO=\boxed{090}$.

 (Solution by TheUltimate123)

 If you're wondering why $KX \cdot KN=KL^2,$ it's because PoP on $(XLN)$ or by $KX \cdot KN=KX \cdot (KX+XN)=KX^2+KX \cdot XN=KX^2+LX^2=KL^2$ (last part by geometric mean theorem / similarity).

 Note: the "semicircle" circumscribing points XOMN is not a semicircle. That is just there to tell you that X, O, M, N are indeed concyclic, so ignore the subtlety of the diagram that makes O seems slightly off the marks than it should be.

 





\textbf{Solution 6 (Alternative PoP)}


                    \begin{figure}[H]
                        \centering
                        \includegraphics[width=0.5\textwidth]{Solution/AIME_2019_I_6_Solution_Figure_1.png}
                    \end{figure}
                    

 (Diagram by vedadehhc)

 Call the base of the altitude from $L$ to $NK$ point $P$. Let $PO=x$. Now, we have that $KP=\sqrt{64-x^2}$ by the Pythagorean Theorem. Once again by Pythagorean, $LO=\sqrt{720+x^2}-x$. Using Power of a Point, we have 

 
 \[(KO)(OM)=(LO)(OQ)\] ($Q$ is the intersection of $OL$ with the circle $\neq L$)

 
 \[8(MO)=(\sqrt{720+x^2}-x)(\sqrt{720+x^2}+x)\]

 
 \[8(MO)=720\]

 
 \[MO=\boxed{090}\].

 (Solution by RootThreeOverTwo)
 \[\]

 \textbf{Remark: Length of OQ} Since $P$ is on the circle’s diameter, $QP = LP = \sqrt{720+x^2}$. So, $OQ = PQ + PO = x + \sqrt{720+x^2}$.~diyarv

 





\textbf{Solution 7 (just one pair of similar triangles)}


                    \begin{figure}[H]
                        \centering
                        \includegraphics[width=0.5\textwidth]{Solution/AIME_2019_I_6_Solution_Figure_1.png}
                    \end{figure}
                    Note that since $\angle KLN = \angle KMN$, quadrilateral $KLMN$ is cyclic. Therefore, we have 
 \[\angle LMK = \angle LNK = 90^{\circ} - \angle LKN = \angle KLP,\]so $\triangle KLO \sim \triangle KML$, giving 
 \[\frac{KM}{28} = \frac{28}{8} \implies KM = 98.\] Therefore, $OM = 98-8 = \boxed{90}$.

 





\textbf{Solution 8 (Pythagoras Bash)}


                    \begin{figure}[H]
                        \centering
                        \includegraphics[width=0.5\textwidth]{Solution/AIME_2019_I_6_Solution_Figure_1.png}
                    \end{figure}
                    

 By Pythagorean Theorem, $KM^2+65^2 = KN^2 = 28^2 + LN^2$. Thus, $LN^2 = KM^2 + 65^2 - 28^2$.

 By Pythagorean Theorem, $KP^2 + LP^2 = 28^2$, and $PN^2 + LP^2 = LN^2$.

 
 \[PN^2 = (KN - KP)^2 = (\sqrt{KM^2 + 65^2} - KP)^2\]

 It follows that
 \[(\sqrt{KM^2 + 65^2} - KP)^2 + LP^2 = KM^2 + 65^2 - 28^2\]

 
 \[KM^2 + 65^2 - 2\sqrt{KM^2 + 65^2}(KP) + KP^2 + LP^2 = KM^2 + 65^2 - 28^2\]

 Since $KP^2 + LP^2 = 28^2$,
 \[KM^2 + 65^2 - 2\sqrt{KM^2 + 65^2}(KP) + 28^2 = KM^2 + 65^2 - 28^2\]

 
 \[-2\sqrt{KM^2 + 65^2}(KP) = -2 \times 28^2\]

 
 \[KP = \frac{28^2}{\sqrt{KM^2 + 65^2}}\]

 $\angle OKP = \angle NKM$ (it's the same angle) and $\angle OPK = \angle KMN = 90^{\circ}$. Thus, $\triangle KOP \sim \triangle KNM$.

 Thus,

 
 \[\frac{KO}{KN} = \frac{KP}{KM}\]

 
 \[\frac{8}{\sqrt{KM^2 + 65^2}} = \frac{\frac{28^2}{\sqrt{KM^2 + 65^2}}}{KM}\]

 Multiplying both sides by $\sqrt{KM^2 + 65^2}$:

 
 \[8 = \frac{28^2}{KM}\]

 
 \[KM = 98\]

 Therefore, $OM = 98-8 = \boxed{90}$

 




	\item (\textit{AIME I 2019}) There are positive integers $x$ and $y$ that satisfy the system of equations

 \[\log_{10} x + 2 \log_{10} (\gcd(x,y)) = 60\] 
 \[\log_{10} y + 2 \log_{10} (\text{lcm}(x,y)) = 570.\]
Let $m$ be the number of (not necessarily distinct) prime factors in the prime factorization of $x$, and let $n$ be the number of (not necessarily distinct) prime factors in the prime factorization of $y$. Find $3m+2n$.


 
 \textbf{Answer: }880



\textbf{Solution 1}

Add the two equations to get that $\log x+\log y+2(\log(\gcd(x,y))+2(\log(\text{lcm}(x,y)))=630$.Then, we use the theorem $\log a+\log b=\log ab$ to get the equation, $\log (xy)+2(\log(\gcd(x,y))+\log(\text{lcm}(x,y)))=630$.Using the theorem that $\gcd(x,y) \cdot \text{lcm}(x,y)=x\cdot y$, along with the previously mentioned theorem, we can get the equation $3\log(xy)=630$.This can easily be simplified to $\log(xy)=210$, or $xy = 10^{210}$.

 $10^{210}$ can be factored into $2^{210} \cdot 5^{210}$, and $m+n$ equals to the sum of the exponents of $2$ and $5$, which is $210+210 = 420$.Multiply by two to get $2m +2n$, which is $840$.Then, use the first equation ($\log x + 2\log(\gcd(x,y)) = 60$) to show that $x$ has to have lower degrees of $2$ and $5$ than $y$  (you can also test when $x>y$, which is a contradiction to the restrains we set before). Therefore, $\gcd(x,y)=x$. Then, turn the equation into $3\log x = 60$, which yields $\log x = 20$, or $x = 10^{20}$.Factor this into $2^{20} \cdot 5^{20}$, and add the two 20's, resulting in $m$, which is $40$.Add $m$ to $2m + 2n$ (which is $840$) to get $40+840 = \boxed{880}$.

 Remark: You can obtain the contradiction by using LTE. If $\nu_2{(x)}\geq{\nu_2{(y)}}, \nu_2{(y^2x)}=60$. However, $\nu_2{(xy)}=210$ a contradiction. Same goes with taking $\nu_5{(x,y)}$

 \textbf{Easier Approach to Finish} After noting that $xy=10^{210},$ notice that we can let $x=10^a$ and $y=10^b.$ Thus, we have from the given equations (1) and (2) respectively, that:
 \[a+2a=60\]
 \[b+2b=570\]Solving, we get $(a,b)=(20,190).$ This matches with our constraint that $xy=10^{210}$ (this constraint can actually be rederived by adding the two equations) so we finish from here.

 $x=2^{20}\cdot 5^{20}, y=2^{190}\cdot 5^{190}.$ Thus, the answer desired is $3(20+20)+2(190+190)=880.$ ~mathboy282 (minor addition by Technodoggo)

 \textbf{Why only 10?} If you didn't understand the main solution 1 or you read the easier finish, a question that might still linger is why can we let  $x=10^a$ and $y=10^b$?

 Notice how if we set $x = 2^a$ and $y = 5^b$, in the first equation we obtain$a\log(2) =60$. However, This implies that $a$ must be a non-integer logarithmic value, which is not allowed as that implies $x$ is not an integer. The same applies for $y$. Thus, we cannot have $x = 2^a$ and $y = 5^b$.

 





\textbf{Solution 2}

First simplifying the first and second equations, we get that 

 
 \[\log_{10}(x\cdot\text{gcd}(x,y)^2)=60\]
 \[\log_{10}(y\cdot\text{lcm}(x,y)^2)=570\]

 
Thus, when the two equations are added, we have that
 \[\log_{10}(x\cdot y\cdot\text{gcd}^2\cdot\text{lcm}^2)=630\]When simplified, this equals 
 \[\log_{10}(x^3y^3)=630\]so this means that
 \[x^3y^3=10^{630}\]so
 \[xy=10^{210}.\]

 Now, the following cannot be done on a proof contest but let's (intuitively) assume that $x<y$ and $x$ and $y$ are both powers of $10$. This means the first equation would simplify to 
 \[x^3=10^{60}\] and 
 \[y^3=10^{570}.\] Therefore, $x=10^{20}$ and $y=10^{190}$ and if we plug these values back, it works! $10^{20}$ has $20\cdot2=40$ total factors and $10^{190}$ has $190\cdot2=380$ so 
 \[3\cdot 40 + 2\cdot 380 = \boxed{880}.\]

 Please remember that you should only assume on these math contests because they are timed; this would technically not be a valid solution.

 \textbf{Proof for WLOG} The following assumption made on solution 2 can be done with a generality proof, and is quite simple.

 If we have \( x = 10^a \) and \( y = 10^b \), we know that \( xy = 10^{a+b} \). Then, because 
 \[xy\] is commutative and so are the powers, it doesn't matter which one of \( a \) and \( b \) is greater. Thus, we may safely assume \( a < b \) (or vice versa), and continue.

 Optionally, in a proof contest, you could also show that both $a < b$ and $a > b$ work as valid cases, and provide the same answer.

 





\textbf{Solution 3 (Easy Solution)}

Let $x=10^a$ and $y=10^b$ and $a<b$. Then the given equations become $3a=60$ and $3b=570$. Therefore, $x=10^{20}=2^{20}\cdot5^{20}$ and $y=10^{190}=2^{190}\cdot5^{190}$. Our answer is $3(20+20)+2(190+190)=\boxed{880}$.

 Note from Wiselion: This solution also uses the assumptions as solution 2. In theory x and y don/t have to be power of 10. Although having it not a power of 10 suggests that the sums won't be integers, if the gcd and lcm logs match up respectively with the remains of  log x and y then the logs could add up to an integer. (We'd have to prove it)(PS never wrote the solution)

 





\textbf{Solution 4}

We will use the notation $(a, b)$ for $\gcd(a, b)$ and $[a, b]$ as $\text{lcm}(a, b)$. We can start with a similar way to Solution 1. We have, by logarithm properties, $\log_{10}{x}+\log_{10}{(x, y)^2}=60$ or $x(x, y)^2=10^{60}$. We can do something similar to the second equation and our two equations become 
 \[x(x, y)^2=10^{60}\] 
 \[y[x, y]^2=10^{570}\]Adding the two equations gives us $xy(x, y)^2[x, y]^2=10^{630}$. Since we know that $(a, b)\cdot[a, b]=ab$, $x^3y^3=10^{630}$, or $xy=10^{210}$. We can express $x$ as $2^a5^b$ and $y$ as $2^c5^d$. Another way to express $(x, y)$ is now $2^{min(a, c)}5^{min(b, d)}$, and $[x, y]$ is now $2^{max(a, c)}5^{max(b, d)}$. We know that $x<y$, and thus, $a<c$, and $b<d$. Our equations for $lcm$ and $gcd$ now become 
 \[2^a5^b(2^a5^a)^2=10^{60}\] or $a=b=20$. Doing the same for the $lcm$ equation, we have $c=d=190$, and $190+20=210$, which satisfies $xy=210$. Thus, $3m+2n=3(20+20)+2(190+190)=\boxed{880}$.~awsomek

 





\textbf{Solution 5}

Let $x=d\alpha, y=d\beta, (\alpha, \beta)=1$. Simplifying, $d^3\alpha=10^{60}, d^3\alpha^2\beta^3=10^{510} \implies \alpha\beta^3 = 10^{510}=2^{510} \cdot 5^{510}$. Notice that since $\alpha, \beta$ are coprime, and $\alpha < 5^{90}$(Prove it yourself !) , $\alpha=1, \beta = 10^{170}$. Hence, $x=10^{20}, y=10^{190}$ giving the answer $\boxed{880}$. 

 (Solution by Prabh1512)

 





\textbf{Solution 6 (Official MAA)}

The two equations are equivalent to $x(\gcd(x,y))^2=10^{60}$ and $y(\operatorname{lcm}(x,y))^2=10^{570}$ respectively. Multiplying corresponding sides of the equations leads to $xy(\gcd(x,y)\operatorname{lcm}(x,y))^2=(xy)^3=10^{630}$, so $xy=10^{210}$. It follows that there are nonnegative integers $a,\,b,\,c,$ and $d$ such that $(x,y)=(2^a5^b,2^c5^d)$ with $a+c=b+d=210$. Furthermore, 
 \[\frac{(\operatorname{lcm}(x,y))^2}{x}=\frac{y(\operatorname{lcm}(x,y))^2}{xy}=\frac{10^{570}}{10^{210}}=10^{360}.\] Thus $\max(2a,2c)-a=\max(2b,2d)-b=360.$ Because neither $2a-a$ nor $2b-b$ can equal $360$ when $a+c=b+d=210,$ it follows that $2c-a=2d-b=360$. Hence $(a,b,c,d)=(20,20,190,190)$, so the prime factorization of $x$ has $20+20=40$ prime factors, and the prime factorization of $y$ has $190+190=380$ prime factors. The requested sum is $3\cdot40+2\cdot380=880.$

 





\textbf{Solution 7}

Add the two equations and use the fact that $\gcd\left(x,y\right)\cdot\mathrm{lcm}\left(x,y\right)=xy$ to find that $xy=10^{210}$. So let $x=2^a5^b$ and $y=2^{210-a}5^{210-b}$ for $0\leq a,b\leq210$. If $a\geq105$ then the exponent of $2$ in $x\cdot\gcd\left(x,y\right)^2=10^{60}$ is $a+2\left(210-a\right)=420-a$, so $a=360$, contradiction. So $a<105$. Then the exponent of $2$ in $x\cdot\gcd\left(x,y\right)^2$ is $a+2a=3a$, so $a=20$. Similarly, $b=20$. Then $3m+2n=3\left(a+b\right)+2\left(420-a-b\right)=\boxed{880}$ as desired.

 





\textbf{Solution 8}

We can simplify the equations step by step. The first equation simplifies to $log($(x)($(gcd(x,y))^2$)$)=60$. The second equation simplifies to log($(y)$($(lcm(x,y)^2$)$)=570$. Up to here, we used the exponent and addition log identities. 

 Now before we move on to the next few simplification steps, we must remember that $gcd(a,b)$*$lcm(a,b)$=$a*b$.

 \begin{verbatim}
Due to the fact that these log's are in base $10$, this makes the first and second equation equal to $10^{60}$, $10^{570}$ respectively. In this step, we switched the log's into exponential form. Now we multiply both equations to get $x*y$*$(xy)^2$=($x^3$)($y^3$)=$10^{630}$. Now we take the cube root of both sides to get $xy=10^{210}$. 
\end{verbatim}We've now gotten to the crucial part of this equation. Though this wouldn't pass for full points in a proof-based contest, this is AIME. So, we assume that $x<y$. We also let $x$=$10^a$ and $y$=$10^b$ That means that $gcd(x,y)$ is $x$ and the $lcm(x,y)$ is $y$ due to the fact that we are also assuming that both $x,y$ are $10^a$, $10^b$ respectively. 

 If we put our last few insights together into the first and second equation, we see that $x$=$10^{20}$. We also see that $y$=$10^{190}$. We could check these if wanted (don't worry they work), but if you were very limited on time for this question, just assume these values work and move on. 

 Now $10^{20}$ factors as $2^{20}$*$5^{20}$. This has $40$ prime factors. $2$, $20$ times and $5$, $20$ times. $10^{190}$ factors as $2^{190}$*$5^{190}$. This has $380$ prime factors. $2$, $190$ times and $5$, $190$ times. Now it's just $40*3+380*2=880$ as our final answer.

 -Schintalpati

 




	\item (\textit{AIME I 2019}) Let $x$ be a real number such that $\sin^{10}x+\cos^{10} x = \frac{11}{36}$. Then $\sin^{12}x+\cos^{12} x = \frac{m}{n}$ where $m$ and $n$ are relatively prime positive integers. Find $m+n$.


 
 \textbf{Answer: }067



\textbf{Solution 1}

We can substitute $y = \sin^2{x}$. Since we know that $\cos^2{x}=1-\sin^2{x}$, we can do some simplification. 

 This yields $y^5+(1-y)^5=\frac{11}{36}$. From this, we can substitute again to get some cancellation through binomials. If we let $z=\frac{1}{2}-y$, we can simplify the equation to:
 \[\left(\frac{1}{2}+z\right)^5+\left(\frac{1}{2}-z\right)^5=\frac{11}{36}.\] After using binomial theorem, this simplifies to:
 \[\frac{1}{16}(80z^4+40z^2+1)=\frac{11}{36}.\]If we use the quadratic formula, we obtain  $z^2=\frac{1}{12}$, so $z=\pm\frac{1}{2\sqrt{3}}$ (observe that either choice of $z$ doesn't matter). Substituting $z,$ we get:

 
 \[\sin^{12}{x}+\cos^{12}{x}=\left(\frac{1}{2}-z\right)^6+\left(\frac{1}{2}+z\right)^6=2z^6 + \frac{15z^4}{2} + \frac{15z^2}{8} + \frac{1}{32}=\frac{13}{54}.\]

 Therefore, the answer is $13+54=\boxed{067}$.

 -eric2020, inspired by Tommy2002

 \textbf{Motivation} The motivation to substitute $z=\frac{1}{2}-y$ comes so that after applying the binomial theorem to $y^5+(1-y)^5=\left(\frac{1}{2}+z\right)^5+\left(\frac{1}{2}-z\right)^5,$ a lot of terms will cancel out. Note that all the terms with odd exponents in $\left(\frac{1}{2}+z\right)^5$ will cancel out, while the terms with even exponents will be doubled. \textbf{mathboy282}

 





\textbf{Solution 2}

First, for simplicity, let $a=\sin{x}$ and $b=\cos{x}$. Note that $a^2+b^2=1$. We then bash the rest of the problem out. Take the fifth power of this expression and get $a^{10}+b^{10}+5a^2b^2(a^6+b^6)+10a^4b^4(a^2+b^2)=\frac{11}{36}+5a^2b^2(a^6+b^6)+10a^4b^4=1$. Note that we also have $\frac{11}{36}=a^{10}+b^{10}=(a^{10}+b^{10})(a^2+b^2)=a^{12}+b^{12}+a^2b^2(a^8+b^8)$. So, it suffices to compute $a^2b^2(a^8+b^8)$. Let $y=a^2b^2$. We have from cubing $a^2+b^2=1$ that $a^6+b^6+3a^2b^2(a^2+b^2)=1$ or $a^6+b^6=1-3y$. Next, using $\frac{11}{36}+5a^2b^2(a^6+b^6)+10a^4b^4=1$, we get $a^2b^2(a^6+b^6)+2a^4b^4=\frac{5}{36}$ or $y(1-3y)+2y^2=y-y^2=\frac{5}{36}$. Solving gives $y=\frac{5}{6}$ or $y=\frac{1}{6}$. Clearly $y=\frac{5}{6}$ is extraneous, so $y=\frac{1}{6}$. Now note that $a^4+b^4=(a^2+b^2)^2-2a^2b^2=\frac{2}{3}$, and $a^8+b^8=(a^4+b^4)^2-2a^4b^4=\frac{4}{9}-\frac{1}{18}=\frac{7}{18}$. Thus we finally get $a^{12}+b^{12}=\frac{11}{36}-\frac{7}{18}\cdot\frac{1}{6}=\frac{13}{54}$, giving $\boxed{067}$.

 \textbf{- Emathmaster}

 





\textbf{Solution 3 (Newton Sums)}

Newton sums is basically constructing the powers of the roots of the polynomials instead of deconstructing them which was done in Solution $2$. Let $\sin^2x$ and $\cos^2x$ be the roots of some polynomial $F(a)$. Then, by Vieta, $F(a)=a^2-a+b$ for some $b=\sin^2x\cdot\cos^2x$.

 Let $S_k=\left(\sin^2x\right)^k+\left(\cos^2x\right)^k$. We want to find $S_6$. Clearly $S_1=1$ and $S_2=1-2b$. Newton sums tells us that $S_k-S_{k-1}+bS_{k-2}=0\Rightarrow S_k=S_{k-1}-bS_{k-2}$ where $k\ge 3$ for our polynomial $F(a)$. 

 Bashing, we have
 \[S_3=S_2-bS_1\Rightarrow S_3=(1-2b)-b(1)=1-3b\]
 \[S_4=S_3-bS_2\Rightarrow S_4=(1-3b)-b(1-2b)=2b^2-4b+1\]
 \[S_5=S_4-bS_3\Rightarrow S_5=(2b^2-4b+1)-b(1-3b)=5b^2-5b+1=\frac{11}{36}\]

 Thus 
 \[5b^2-5b+1=\frac{11}{36}\Rightarrow 5b^2-5b+\frac{25}{36}=0, 36b^2-36b+5=0, (6b-1)(6b-5)=0\]$b=\frac{1}{6} \text{ or } \frac{5}{6}$. Clearly, $\sin^2x\cdot\cos^2x\not=\frac{5}{6}$ so $\sin^2x\cdot\cos^2x=b=\frac{1}{6}$.

 Note $S_4=\frac{7}{18}$. Solving for $S_6$, we get $S_6=S_5-\frac{1}{6}S_4=\frac{13}{54}$. Finally, $13+54=\boxed{067}$.

 





\textbf{Solution 4}

Factor the first equation. 
 \[\sin^{10}x + \cos^{10}x = (\sin^2x+\cos^2x)(\sin^8x-\sin^6x\cos^2x+\sin^4x\cos^4x-\sin^2x\cos^6x+\cos^8x)\]First of all, $\sin^4x+\cos^4x = 1-2\sin^2x\cos^2x$ because $\sin^4x+\cos^4x=(\sin^2x + \cos^2x)^2 -2\sin^2x\cos^2x = 1 - 2\sin^2x\cos^2x$We group the first, third, and fifth term and second and fourth term. The first group: 
 \begin{align*} \sin^8+\sin^4x\cos^4x+\cos^8x &= (\sin^4x+\cos^4x)^2-\sin^4x\cos^4x)\\ &= (1 - 2\sin^2x\cos^2x)^2-\sin^4x\cos^4x)\\ &= 1+4\sin^4x\cos^4x-4\sin^2x\cos^2x \end{align*}
The second group: 
 \begin{align*} -\sin^6x\cos^2x-\sin^2x\cos^6x  &= -\sin^2x\cos^2x(\sin^4x+\cos^4x)\\ &= -\sin^2x\cos^2x(1-2\sin^2x\cos^2x)\\ &= -\sin^2x\cos^2x+2\sin^4x\cos^4x \end{align*}
Add the two together to make 
 \[1+4\sin^4x\cos^4x-4\sin^2x\cos^2x-\sin^2x\cos^2x+2\sin^4x\cos^4x = 1 - 5\sin^2x\cos^2x+5\sin^4x\cos^4x\] Because this equals $\frac{11}{36}$, we have 
 \[5\sin^4x\cos^4x- 5\sin^2x\cos^2x+\frac{25}{36}=0\] Let $\sin^2x\cos^2x = a$ so we get 
 \[5a^2- 5a+\frac{25}{36}=0 \Rightarrow a^2-a+\frac{5}{36}\] Solving the quadratic gives us 
 \[a = \frac{1 \pm \frac{2}{3}}{2}\] Because $\sin^2x\cos^2x \le \frac{1}{4}$, we finally get $a = \frac{1 - \frac{2}{3}}{2} = \frac{1}{6}$. 

 Now from the second equation, 
 \begin{align*} \sin^{12}x + \cos^{12}x &= (\sin^4x+\cos^4x)(\sin^8x-\sin^4x\cos^4x+\cos^8x)\\ &= (1-2\sin^2x\cos^2x)((\sin^4x+\cos^4x)^2-3\sin^4x\cos^4x)\\ &= (1-2\sin^2x\cos^2x)((1-2\sin^2x\cos^2x)^2-3\sin^4x\cos^4x) \end{align*}
Plug in $\sin^2x\cos^2x = \frac{1}{6}$ to get 
 \[(1-2(\frac{1}{6}))(1-2(\frac{1}{6})^2-3(\frac{1}{6})^2) = \frac{13}{54}\]which yields the answer $\boxed{067}$

 





\textbf{Solution 5}

Define the recursion $a_n=(\sin^2 x)^n+(\cos^2 x)^n$We know that the characteristic equation of $a_n$ must have 2 roots, so we can recursively define $a_n$ as $a_n=p*a_{n-1}+q*a_{n-2}$. $p$ is simply the sum of the roots of the characteristic equation, which is $\sin^2 x+\cos^2 x=1$. $q$ is the product of the roots, which is $-(\sin^2 x)(\cos^2 x)$. This value is not trivial and we have to solve for it.We know that $a_0=2$, $a_1=1$, $a_5=\frac{11}{36}$.Solving the rest of the recursion gives

 
 \[a_2=1+2q\]
 \[a_3=1+3q\]
 \[a_4=1+4q+2q^2\]
 \[a_5=1+5q+5q^2=\frac{11}{36}\]
 \[a_6=1+6q+9q^2+2q^3\]

 
Solving for $q$ in the expression for $a_5$ gives us $q^2+q+\frac{5}{36}=0$, so $q=-\frac{5}{6}, -\frac{1}{6}$. Since $q=-(\sin^2 x)(\cos^2 x)$, we know that the minimum value it can attain is $-\frac{1}{4}$ by AM-GM, so $q$ cannot be $-\frac{5}{6}$.Plugging in the value of $q$ into the expression for $a_6$, we get $a_6=1-1+\frac{1}{4}-\frac{1}{108}=\frac{26}{108}=\frac{13}{54}$. Our final answer is then $13+54=\boxed{067}$

 -Natmath

 





\textbf{Solution 6}

Let $m=\sin^2 x$ and $n=\cos^2 x$, then $m+n=1$ and $m^5+n^5=\frac{11}{36}$

 $m^6+n^6=(m^5+n^5)(m+n)-mn(m^4+n^4)=(m^5+n^5)-mn(m^4+n^4)$

 Now factoring $m^5+n^5$ as solution 4 yields $m^5+n^5=(m+n)(m^4-m^3n+m^2n^2-mn^3+n^4)$$=m^4+n^4-mn(m^2-mn+n^2)=m^4+n^4-mn[(m+n)^2-3mn]=m^4+n^4-mn(1-3mn)$. 

 Since $(m+n)^4=m^4+4m^3n+6m^2n^2+4mn^3+n^4$, $m^4+n^4=(m+n)^4-2mn(2m^2+3mn+2n^2)=1-2mn(2m^2+3mn+2n^2)$. 

 Notice that $2m^2+3mn+2n^2$ can be rewritten as $[\sqrt{2}(a+b)]^2-mn=2-mn$. Thus,$m^4+n^4=1-2mn(2-mn)$ and $m^5+n^5=1-2mn(2-mn)-mn(1-3mn)=1-5mn+5(mn)^2=\frac{11}{36}$. As in solution 4, we get $mn=\frac{1}{6}$ and $m^4+n^4=1-2*\frac{1}{6}(2-\frac{1}{6})=\frac{7}{18}$

 Substitute $m^4+n^4=\frac{7}{18}$ and $mn=\frac{1}{6}$, then$m^6+n^6=\frac{11}{36}-\frac{1}{6}*\frac{7}{18}=\frac{13}{54}$, and the desired answer is $\boxed{067}$

 





\textbf{Solution 7 (Algebra and Recursion)}

This was my solution on the real test. Let $a=\sin^{2}x, b=\cos^{2}x$. Observe that $a+b=1$, and $(a^{n-1}+b^{n-1}) \cdot (a+b) = a^{n-1}b+ab^{n-1}+a^n+b^n=a^{n-1}+b^{n-1} \implies a^n+b^n=a^{n-1}+b^{n-1}-ab(a^{n-2}+b^{n-2})$. Let $ab=x$, and we want to know $x$. 

 Starting from $n=2$, we have $a^2+b^2=1-2x \implies a^3+b^3=1-2x-x=1-3x \implies a^4+b^4=1-4x-x(1-2x)=1-5x+2x^2 \implies a^5+b^5=2x^2-5x+1-x(1-3x)=5x^2-5x+1=\frac{11}{36}$. Rearrange this quadratic and divide by $5$ to get $x^2-x+\frac{5}{36}=0 \implies x=\{\frac{1}{6}, \frac{5}{6}\}$. (You can guess the factorization or use the quadratic formula.) 

 Given we have two solutions, plug both in and see which results in a positive rational fraction. Plugging in $x=\frac{1}{6}$, we get $a^2+b^2=\frac{2}{3}, a^3+b^3=\frac{1}{2}, a^4+b^4=\frac{7}{18}, a^5+b^5=\frac{11}{36}, a^6+b^6=\frac{11}{36}-\frac{1}{6} \cdot \frac{7}{18} = \frac{26}{108} = \frac{13}{54} \implies \boxed{67}$. 

 





\textbf{Solution 8 (Official MAA)}

Let $c=\sin^2x\cdot\cos^2x,$ and let $S(n)=\sin^{2n}x+\cos^{2n}x.$ Then for $n\ge 1$
 \begin{align*} S(n)&=(\sin^{2n}x+\cos^{2n}x)\cdot(\sin^2x+\cos^2x)\\ &=\sin^{2n+2}x+\cos^{2n+2}x+\sin^2x\cdot\cos^2x(\sin^{2n-2}x+\cos^{2n-2}x)\\ &=S(n+1)+cS(n-1). \end{align*}
Because $S(0)=2$ and $S(1)=1,$ it follows that $S(2)=1-2c, S(3)=1-3c,S(4)=2c^2-4c+1,$ and $\frac{11}{36}=S(5)=5c^2-5c+1.$ Hence $c=\frac16$ or $\frac56,$ and because $4c=\sin^2{2x},$ the only possible value of $c$ is $\frac16.$ Therefore 
 \[S(6)=S(5)-cS(4)=\frac{11}{36}-\frac16\left(2\left(\frac16\right)^2-4\left(\frac16\right)+1\right)=\frac{13}{54}.\] The requested sum is $13+54=67.$

 





\textbf{Solution 9 (Recursion)}

Let $a_n=\sin^nx+\cos^nx$ for non-negative integers $n$. Then $a_0=2$ and $a_2=1$. In addition,
 \[a_n=\sin^nx+\cos^nx=\left(\sin^{n-2}x+\cos^{n-2}x\right)\left(\sin^2x+\cos^2x\right)-\sin^2x\cos^2x\left(\sin^{n-4}x+\cos^{n-4}x\right)=a_{n-2}-Xa_{n-4},\]where $X=\sin^2x\cos^2x$. So we can compute\begin{align*}a_4&=1-2X\\a_6&=1-3X\\a_8&=1-4X+2X^2\\a_{10}&=1-5X+5X^2=\frac{11}{36}\end{align*}
so $X=\frac{1}{6},\frac{5}{6}$. But by the sin double angle formula, $\sin^2x\cos^2x=\frac{1}{4}\sin^22x\leq\frac{1}{4}$, so $X=\frac{1}{6}$. Then
 \[a_{12}=a_{10}-Xa_8=\frac{11}{36}-\frac{1}{6}\cdot\frac{7}{18}=\frac{13}{54}\]so the answer is $\boxed{067}$ as desired.

 A quick note: this solution uses recursion and is similar to the solution 7 above. It was from trumpeter, posted in the AoPS Forums, Contest Discussion.

 





\textbf{Solution 10 (Algebra but not bad)}

Note that from the Pythagorean identity, $\sin^{2}(x)+\cos^{2}(x)=1$. Now, let $m=\sin^{2}(x)$ and let $n=\cos^{2}(x)$. Now, we seek $m^{6}+n^{6}$ given

 $m+n=1$, and$m^{5}+n^{5}=\frac{11}{36}$. 

 Now, note the factorization $m^{5}+n^{5}=(m+n)(m^{4}-m^{3}n+m^{2}n^{2}-mn^{3}+n^{4})=\frac{11}{36}$, so
 \[m^{4}-m^{3}n+m^{2}n^{2}-mn^{3}+n^{4}=\frac{11}{36}\].

 Now, we also have $(m+n)^{4}=m^{4}+4m^{3}n+6m^{2}n^{2}+4mn^{2}+n^{4}=1$. Now, upon subtraction, we obtain
 \[5(m^{3}n+m^{2}n^{2}+mn^{3})=\frac{25}{36} \rightarrow m^{3}n+m^{2}n^{2}+mn^{3}=\frac{5}{36}\].

 This tells us that $mn(m^{2}+mn+n^{2})=\frac{5}{36}$. However, note that $m^{2}+2mn+n^{2}=(m+n)^{2}=1$, so if we let $a=mn$, we have
 \[a(1-a)=\frac{5}{36}\].

 Solving this quadratic yields $a=\frac{1}{6}$ and $a=\frac{5}{6}$. Now, note that $a \neq \frac{5}{6}$ because of the fact that $\sin^{2}(x) \cos^{2}(x)$ is really small, so $a=\frac{1}{6}$. 

 Now, we seek $m^{6}+n^{6}$. This can be found by finding $(m^{3}+n^{3})^{2}-2m^{3}n^{3}$. 

 We have $m^{3}+n^{3}=(m+n)(m^{2}-mn+n^{2})=(1)\left(1^{2}-3\left(\frac{1}{6}\right)\right)=\frac{1}{2}$. Thus, our desired value is

 
 \[\left(\frac{1}{2}\right)^{2}-2\left(\frac{1}{6}\right)^{3}=\frac{1}{4}-\frac{1}{108}=\frac{13}{54}\]and the desired extraction is $13+54=\fbox{067}$

 




	\item (\textit{AIME I 2019}) Let $\tau(n)$ denote the number of positive integer divisors of $n$. Find the sum of the six least positive integers $n$ that are solutions to $\tau (n) + \tau (n+1) = 7$.


 
 \textbf{Answer: }540



\textbf{Solution}

In order to obtain a sum of $7$, we must have:

 
                \begin{itemize}
                    \item either a number with $5$ divisors (\textit{a fourth power of a prime}) and a number with $2$ divisors (\textit{a prime}), or\item a number with $4$ divisors (\textit{a semiprime or a cube of a prime}) and a number with $3$ divisors (\textit{a square of a prime}). (No integer greater than $1$ can have fewer than $2$ divisors.)
                \end{itemize}


                Since both of these cases contain a number with an odd number of divisors, that number must be an even power of a prime. These can come in the form of a square-like $3^2$ with $3$ divisors, or a fourth power like $2^4$ with $5$ divisors. We then find the smallest such values by hand.

 
                \begin{itemize}
                    \item $2^2$ has two possibilities: $3$ and $4$ or $4$ and $5$. Neither works.\item $3^2$ has two possibilities: $8$ and $9$ or $9$ and $10$. $\boxed{(8,9)}$ and $\boxed{(9,10)}$ both work.\item $2^4$ has two possibilities: $15$ and 
                \end{itemize}


                




	\item (\textit{AIME I 2019}) For distinct complex numbers $z_1,z_2,\dots,z_{673}$, the polynomial 

 \[(x-z_1)^3(x-z_2)^3 \cdots (x-z_{673})^3\]
can be expressed as $x^{2019} + 20x^{2018} + 19x^{2017}+g(x)$, where $g(x)$ is a polynomial with complex coefficients and with degree at most $2016$. The value of 

 \[\left| \sum_{1 \le j <k \le 673} z_jz_k \right|\]
can be expressed in the form $\frac{m}{n}$, where $m$ and $n$ are relatively prime positive integers. Find $m+n$.


 
 \textbf{Answer: }352



\textbf{Solution 1}

In order to begin this problem, we must first understand what it is asking for. The notation 
 \[\left| \sum_{1 \le j <k \le 673} z_jz_k \right|\]simply asks for the absolute value of the sum of the product of the distinct unique roots of the polynomial taken two at a time or
 \[(z_1z_2+z_1z_3+ \dots + z_1z_{672}+z_1z_{673})+(z_2z_3+z_2z_4+ \dots +z_2z_{673}) + (z_3z_4+z_3z_5+ \dots +z_3z_{673}) + \dots +z_{672}z_{673}.\] Call this sum $S$. 

 Now we can begin the problem. Rewrite the polynomial as $P=(x-z_1)(x-z_1)(x-z_1)(x-z_2)(x-z_2)(x-z_2) \dots (x-z_{673})(x-z_{673})(x-z_{673})$. Then we have that the roots of $P$ are $z_1,z_1,z_1,z_2,z_2,z_2, \dots , z_{673},z_{673},z_{673}$. 

 By Vieta's formulas, we have that the sum of the roots of $P$ is $(-1)^1 \cdot \dfrac{20}{1}=-20=z_1+z_1+z_1+z_2+z_2+z_2+ \dots + z_{673}+z_{673}+z_{673}=3(z_1+z_2+z_3+ \dots +z_{673})$. Thus, $z_1+z_2+z_3+ \dots +z_{673}=- \dfrac{20}{3}.$ 

 Similarly, we also have that the the sum of the roots of $P$ taken two at a time is $(-1)^2 \cdot \dfrac{19}{1} = 19.$ This is equal to $z_1^2+z_1^2+z_1^2+z_1z_2+z_1z_2+z_1z_2+ \dots =  \\ 3(z_1^2+z_2^2+ \dots + z_{673}^2) + 9(z_1z_2+z_1z_3+z_1z_4+ \dots + z_{672}z_{673}) =  3(z_1^2+z_2^2+ \dots + z_{673}^2) + 9S.$ 

 Now we need to find and expression for $z_1^2+z_2^2+ \dots + z_{673}^2$ in terms of $S$. We note that $(z_1+z_2+z_3+ \dots +z_{673})^2= (-20/3)^2=\dfrac{400}{9} \\ =(z_1^2+z_2^2+ \dots + z_{673}^2)+2(z_1z_2+z_1z_3+z_1z_4+ \dots + z_{672}z_{673})=(z_1^2+z_2^2+ \dots + z_{673}^2)+2S.$ Thus, $z_1^2+z_2^2+ \dots + z_{673}^2= \dfrac{400}{9} -2S$. 

 Plugging this into our other Vieta equation, we have $3 \left( \dfrac{400}{9} -2S \right) +9S = 19$. This gives $S = - \dfrac{343}{9} \Rightarrow \left| S \right| = \dfrac{343}{9}$. Since 343 is relatively prime to 9, $m+n = 343+9 = \fbox{352}$.

 





\textbf{Solution 2 (Fake Solve)}

This is a quick fake solve using $z_i = 0$ where $3 \le i \le 673$ and only $z_1,z_2 \neq 0$ .

 By Vieta's, 
 \[3z_1+3z_2=-20\] and 
 \[3z_1^2+3z_2^2+9z_1z_2 = 19.\]Rearranging gives $z_1 + z_2 = \dfrac{-20}{3}$ and $3(z_1^2 + 2z_1z_2 + z_2^2) + 3z_1z_2 = 19$ giving $3(z_1 + z_2)^2 + 3z_1z_2 = 19$.

 Substituting gives $3\left(\dfrac{400}{9}\right) + 3z_1z_2 = 19$ which simplifies to $\dfrac{400}{3} + 3z_1z_2 = \dfrac{57}{3}$.

 So, $3z_1z_2 = \dfrac{-343}{3}$, $z_1z_2 = \dfrac{-343}{9}$, $|\dfrac{-343}{9}|=\dfrac{343}{9}$, 
 \[m+n = 343+9 = \boxed{352}.\]

 





\textbf{Solution 3}

Let $x=\sum_{1\le j<k\le 673} z_jz_k$. By Vieta's, 
 \[3\sum_{i=1}^{673}z_i=-20\implies \sum_{i=1}^{673}z_i=-\frac{20}3.\]Then, consider the $19x^{2017}$ term. To produce the product of two roots, the two roots can either be either $(z_i,z_i)$ for some $i$, or $(z_j,z_k)$ for some $j<k$. In the former case, this can happen in $\tbinom 32=3$ ways, and in the latter case, this can happen in $3^2=9$ ways. Hence,
 \begin{align*} 19=3\sum_{i=1}^{673} z_i^2+9\sum_{1\le j<k\le 673} z_jz_k=3\left(\left(-\frac{20}3\right)^2-2x\right)+9x&=\frac{400}3+3x\\ \implies x&=-\frac{343}9, \end{align*}
and the requested sum is $343+9=\boxed{352}$.

 





\textbf{Solution 4}

Let 
 \[(f(x))^3=x^{2019}+20x^{2018}+19x^{2017}+g(x).\] Therefore, $f(x)=(x-z_{1})(x-z_{1})\cdots (x-z_{673})$. This is also equivalent to 
 \[f(x)=x^{673}+ax^{672}+bx^{671}+h(x)\] for some real coefficients $a$ and $b$ and some polynomial $h(x)$ with degree $670$. We can see that the big summation expression is simply summing the product of the roots of $f(x)$ taken two at a time. By Vieta's, this is just the coefficient $b$. The first three terms of $(f(x))^3$ can be bashed in terms of $a$ and $b$ to get 
 \[20 = 3a\]
 \[19 = 3a^2+3b\]Thus, $a=\frac{20}{3}$ and $b=\frac{1}{3}\left(19-3\left(\frac{20}{3} \right)^2 \right)$. That is $|b|=\frac{343}{9}=\frac{m}{n}$. $m+n=343+9=\boxed{352}$

 





\textbf{Solution 5 (Newton's Sums)}

The polynomial has 673 unique roots of multiplicity 3 each, for a total of 2019 roots.

 To start, let $S$ be $\left| \sum_{1 \le j <k \le 673} z_jz_k \right|$. Next, let the sum of the roots of the polynomial be $P_1$, which is equivalent to the expression $3(z_1+z_2+\cdots+z_{673})$, and the sum of squares of these roots be $P_2$, equivalent to $3(z_1^2+z_2^2+z_3^2+\dots+z_{673}^2)$.

 By Vieta's,
 \[3( z_1 + z_2 + z_3 + z_4 \dots + z_{673}) = -20\]
 \[3( z_1^2 + z_2^2 + z_3^2 + \dots + z_{673}^2) + 9S = 19\]

 Therefore, $P_1 = -20$. Using the definition of $P_2$ provided above, \href{https://artofproblemsolving.com/wiki/index.php?title=Newton\%27s\_Sums}{ Newton Sums} tells us that $P_2 + -20 \cdot P_1 \cdot + 19 \cdot 2$

 This means that $P_2 = 362$. Plugging back into the second Vieta's equation yields
 \[362 + 9S = 19 \implies S = -\frac{343}{9}\]Taking the absolute value, our answer is therefore $343 + 9 = \boxed{352}$

 
Note: The tripled roots might seem confusing, but it's easily verified that our Vieta's system of equations is correct using combinatorics to find out how many terms we should have. The right hand side of the first equation has $3\cdot673 = 2019$ terms (exactly what we expect). Unsimplified, $P_2$ also has $2019$ terms and $S$ has $\binom{673}{2}$ terms. We want a total of $\binom{2019}{2}$ terms, which happens to be equal to $2019 + 9\cdot\binom{673}{2}$ so this is correct. 

 





\textbf{Solution 6 (Official MAA 1)}

Because each root of the polynomial appears with multiplicity $3,$ Viète's Formulas show that 
 \[z_1+z_2+\cdots+z_{673}=-\frac{20}3\] and 
 \[z_1^2+z_2^2+\cdots+z_{673}^2 =\frac13\left((-20)^2-2\cdot19\right)=\frac{362}3.\] Then the identity 
 \[\left(\sum_{i=1}^{673}z_i\right)^2=\sum_{i=1}^{673}z_i^2+2\left(\sum_{1\le j<k\le 673}z_jz_k\right)\] shows that 
 \[\sum_{1\le j<k\le 673}z_jz_k=\frac{\left(-\frac{20}3\right)^2-\frac{362}3}{2}=-\frac{343}9.\] The requested sum is $343+9=352.$

 Note that such a polynomial does exist. For example, let $z_{673}=-\frac{20}3,$ and for $i=1,2,3,\dots,336,$ let 
 \[z_i=\sqrt{\frac{343i}{9\sum_{j=1}^{336}j}}\qquad \text{and}\qquad z_{i+336}=-z_i.\] Then 
 \[\sum_{i=1}^{673}z_i=-\frac{20}3\qquad\text{and}\qquad\sum_{i=1}^{673}z_i^2=2\sum_{i=1}^{336}\frac{343i}{9\sum_{i=1}^{336}j}+\left(\frac{20}3\right)^2=\frac{362}3,\] as required.

 





\textbf{Solution 7 (Official MAA 2)}

There are constants $a$ and $b$ such that 
 \[(x-z_1)(x-z_2)(x-z_3)\cdots(x-z_{673})=x^{673}+ax^{672}+bx^{671}+\cdots.\] Then 
 \[(x^{673}+ax^{672}+bx^{671}+\cdots)^3=x^{2019}+20x^{2018}+19x^{2017}+\cdots.\] Comparing the $x^{2018}$ and $x^{2017}$ coefficients shows that $3a=20$ and $3a^2+3b=19.$ Solving this system yields $a=\frac{20}3$ and $b=-\frac{343}9.$ Viète's Formulas then give $\left|\sum_{1\le j<k\le 673}z_jz_k\right|=|b|=\frac{343}9,$ as above.

 





\textbf{Solution 8}

Note that we are trying to find the sum of the products of the pairwise distinct $z_k$. First, disregard the repeated roots and call the roots of the polynomial $x_1, x_2\ldots x_{2019}$ for simplicity's sake. 

 Then, $(z_1+z_2\ldots z_{2019})^2=2\sum_{1\leq j<k\leq2019}{z_jz_k}+\sum_{1\leq k\leq 2019}{z_k^2}$. Notice that $\sum_{1\leq j<k\leq2019} {z_jz_k}=19$ and $(x_1+ x_2\ldots +x_{2019})^2=20^2=400$ by Vietas so $\sum_{1\leq k\leq 2019}{z_k^2}=362$. Also, since every root is repeated 3 times, $\sum_{1\leq k\leq 2019}{z_k^2}=3\sum_{1\leq k\leq673}{z_k^2} \Longrightarrow \sum_{1\leq k\leq673}{z_k^2}=\frac{362}{3}$ where the $z_k$ in the second sum are all distinct.

 Call $|\sum_{1 \le j <k \le 673} z_jz_k|=|s|$. Taking into account repeated roots, if we pair any two roots together, there will be $\binom{3}{1}=3$ copies for each $z_k^2, 1\leq z\leq 673$ and $3\cdot 3=9$ copies for each $z_jz_k$ where $z_j$ and $z_k$ are distinct.

 Since the sum of any two pairs of roots is 19, $9s+3\sum_{1\leq k\leq673}{z_k^2}=19$ or $9s+362=19 \Longrightarrow s=\frac{-343}{9}, |s|=\frac{343}{9}$. $m+n=343+9=352$.

 




	\item (\textit{AIME I 2019}) In $\triangle ABC$, the sides have integer lengths and $AB=AC$. Circle $\omega$ has its center at the incenter of $\triangle ABC$. An \textit{excircle} of $\triangle ABC$ is a circle in the exterior of $\triangle ABC$ that is tangent to one side of the triangle and tangent to the extensions of the other two sides. Suppose that the excircle tangent to $\overline{BC}$ is internally tangent to $\omega$, and the other two excircles are both externally tangent to $\omega$. Find the minimum possible value of the perimeter of $\triangle ABC$.


 
 \textbf{Answer: }020



\textbf{Solution 1}

Let the tangent circle be $\omega$. Some notation first: let $BC=a$, $AB=b$, $s$ be the semiperimeter, $\theta=\angle ABC$, and $r$ be the inradius. Intuition tells us that the radius of $\omega$ is $r+\frac{2rs}{s-a}$ (using the exradius formula). However, the sum of the radius of $\omega$ and $\frac{rs}{s-b}$ is equivalent to the distance between the incenter and the the $B/C$ excenter. Denote the B excenter as $I_B$ and the incenter as $I$. Lemma: $I_BI=\frac{2b*IB}{a}$We draw the circumcircle of $\triangle ABC$. Let the angle bisector of $\angle ABC$ hit the circumcircle at a second point $M$. By the incenter-excenter lemma, $AM=CM=IM$. Let this distance be $\alpha$. Ptolemy's theorem on $ABCM$ gives us 
 \[a\alpha+b\alpha=b(\alpha+IB)\to \alpha=\frac{b*IB}{a}\] Again, by the incenter-excenter lemma, $II_B=2IM$ so $II_b=\frac{2b*IB}{a}$ as desired.Using this gives us the following equation: 
 \[\frac{2b*IB}{a}=r+\frac{2rs}{s-a}+\frac{rs}{s-b}\] Motivated by the $s-a$ and $s-b$, we make the following substitution: $x=s-a, y=s-b$This changes things quite a bit. Here's what we can get from it: 
 \[a=2y, b=x+y, s=x+2y\] It is known (easily proved with Heron's and $a=rs$) that 
 \[r=\sqrt{\frac{(s-a)(s-b)(s-b)}{s}}=\sqrt{\frac{xy^2}{x+2y}}\] Using this, we can also find $IB$: let the midpoint of $BC$ be $N$. Using Pythagorean's Theorem on $\triangle INB$, 
 \[IB^2=r^2+(\frac{a}{2})^2=\frac{xy^2}{x+2y}+y^2=\frac{2xy^2+2y^3}{x+2y}=\frac{2y^2(x+y)}{x+2y}\] We now look at the RHS of the main equation: 
 \[r+\frac{2rs}{s-a}+\frac{rs}{s-b}=r(1+\frac{2(x+2y)}{x}+\frac{x+2y}{y})=r(\frac{x^2+5xy+4y^2}{xy})=\frac{r(x+4y)(x+y)}{xy}=\frac{2(x+y)IB}{2y}\]Cancelling some terms, we have 
 \[\frac{r(x+4y)}{x}=IB\] Squaring, 
 \[\frac{2y^2(x+y)}{x+2y}=\frac{(x+4y)^2*xy^2}{x^2(x+2y)}\to \frac{(x+4y)^2}{x}=2(x+y)\] Expanding and moving terms around gives 
 \[(x-8y)(x+2y)=0\to x=8y\] Reverse substituting, 
 \[s-a=8s-8b\to b=\frac{9}{2}a\] Clearly the smallest solution is $a=2$ and $b=9$, so our answer is $2+9+9=\boxed{020}$-franchester

 





\textbf{Solution 2 (Lots of Pythagorean Theorem)}


                    \begin{figure}[H]
                        \centering
                        \includegraphics[width=0.5\textwidth]{Solution/AIME_2019_I_11_Solution_Figure_1.png}
                    \end{figure}
                    

 First, assume $BC=2$ and $AB=AC=x$. The triangle can be scaled later if necessary. Let $I$ be the incenter and let $r$ be the inradius. Let the points at which the incircle intersects $AB$, $BC$, and $CA$ be denoted $M$, $N$, and $O$, respectively.

 
Next, we calculate $r$ in terms of $x$. Note the right triangle formed by $A$, $I$, and $M$. The length $IM$ is equal to $r$. Using the Pythagorean Theorem, the length $AN$ is $\sqrt{x^2-1}$, so the length $AI$ is $\sqrt{x^2-1}-r$. Note that $BN$ is half of $BC=2$, and by symmetry caused by the incircle, $BN=BM$ and $BM=1$, so $MA=x-1$. Applying the Pythagorean Theorem to $AIM$, we get
 \[r^2+(x-1)^2=\left(\sqrt{x^2-1}-r\right)^2.\]Expanding yields
 \[r^2+x^2-2x+1=x^2-1-2r\sqrt{x^2-1}+r^2,\]which can be simplified to
 \[2r\sqrt{x^2-1}=2x-2.\]Dividing by $2$ and then squaring results in
 \[r^2(x^2-1)=(x-1)^2,\]and isolating $r^2$ gets us
 \[r^2=\frac{(x-1)^2}{x^2-1}=\frac{(x-1)^2}{(x+1)(x-1)}=\frac{x-1}{x+1},\]so $r=\sqrt{\frac{x-1}{x+1}}$.

 
We then calculate the radius of the excircle tangent to $BC$. We denote the center of the excircle $E_N$ and the radius $r_N$.

 Consider the quadrilateral formed by $M$, $I$, $E_N$, and the point at which the excircle intersects the extension of $AB$, which we denote $H$. By symmetry caused by the excircle, $BN=BH$, so $BH=1$.

 Note that triangles $MBI$ and $NBI$ are congruent, and $HBE$ and $NBE$ are also congruent. Denoting the measure of angles $MBI$ and $NBI$ measure $\alpha$ and the measure of angles $HBE$ and $NBE$ measure $\beta$, straight angle $MBH=2\alpha+2\beta$, so $\alpha + \beta=90^\circ$. This means that angle $IBE$ is a right angle, so it forms a right triangle.

 Setting the base of the right triangle to $IE$, the height is $BN=1$ and the base consists of $IN=r$ and $EN=r_N$. Triangles $INB$ and $BNE$ are similar to $IBE$, so $\frac{IN}{BN}=\frac{BN}{EN}$, or $\frac{r}{1}=\frac{1}{r_N}$. This makes $r_N$ the reciprocal of $r$, so $r_N=\sqrt{\frac{x+1}{x-1}}$.

 
Circle $\omega$'s radius can be expressed by the distance from the incenter $I$ to the bottom of the excircle with center $E_N$. This length is equal to $r+2r_N$, or $\sqrt{\frac{x-1}{x+1}}+2\sqrt{\frac{x+1}{x-1}}$. Denote this value $r_\omega$.

 
Finally, we calculate the distance from the incenter $I$ to the closest point on the excircle tangent to $AB$, which forms another radius of circle $\omega$ and is equal to $r_\omega$. We denote the center of the excircle $E_M$ and the radius $r_M$. We also denote the points where the excircle intersects $AB$ and the extension of $BC$ using $J$ and $K$, respectively. In order to calculate the distance, we must find the distance between $I$ and $E_M$ and subtract off the radius $r_M$.

 We first must calculate the radius of the excircle. Because the excircle is tangent to both $AB$ and the extension of $AC$, its center must lie on the angle bisector formed by the two lines, which is parallel to $BC$. This means that the distance from $E_M$ to $K$ is equal to the length of $AN$, so the radius is also $\sqrt{x^2-1}$.

 Next, we find the length of $IE_M$. We can do this by forming the right triangle $IAE_M$. The length of leg $AI$ is equal to $AN$ minus $r$, or $\sqrt{x^2-1}-\sqrt{\frac{x-1}{x+1}}$. In order to calculate the length of leg $AE_M$, note that right triangles $AJE_M$ and $BNA$ are congruent, as $JE_M$ and $NA$ share a length of $\sqrt{x^2-1}$, and angles $E_MAJ$ and $NAB$ add up to the right angle $NAE_M$. This means that $AE_M=BA=x$.

 Using Pythagorean Theorem, we get
 \[IE_M=\sqrt{\left(\sqrt{x^2-1}-\sqrt{\frac{x-1}{x+1}}\right)^2+x^2}.\]Bringing back
 \[r_\omega=IE_M-r_M\]and substituting in some values, the equation becomes
 \[r_\omega=\sqrt{\left(\sqrt{x^2-1}-\sqrt{\frac{x-1}{x+1}}\right)^2+x^2}-\sqrt{x^2-1}.\]Rearranging and squaring both sides gets
 \[\left(r_\omega+\sqrt{x^2-1}\right)^2=\left(\sqrt{x^2-1}-\sqrt{\frac{x-1}{x+1}}\right)^2+x^2.\]Distributing both sides yields
 \[r_\omega^2+2r_\omega\sqrt{x^2-1}+x^2-1=x^2-1-2\sqrt{x^2-1}\sqrt{\frac{x-1}{x+1}}+\frac{x-1}{x+1}+x^2.\]Canceling terms results in
 \[r_\omega^2+2r_\omega\sqrt{x^2-1}=-2\sqrt{x^2-1}\sqrt{\frac{x-1}{x+1}}+\frac{x-1}{x+1}+x^2.\]Since
 \[-2\sqrt{x^2-1}\sqrt{\frac{x-1}{x+1}}=-2\sqrt{(x+1)(x-1)\frac{x-1}{x+1}}=-2(x-1),\]We can further simplify to
 \[r_\omega^2+2r_\omega\sqrt{x^2-1}=-2(x-1)+\frac{x-1}{x+1}+x^2.\]Substituting out $r_\omega$ gets
 \[\left(\sqrt{\frac{x-1}{x+1}}+2\sqrt{\frac{x+1}{x-1}}\right)^2+2\left(\sqrt{\frac{x-1}{x+1}}+2\sqrt{\frac{x+1}{x-1}}\right)\sqrt{x^2-1}=-2(x-1)+\frac{x-1}{x+1}+x^2\]which when distributed yields
 \[\frac{x-1}{x+1}+4+4\left(\frac{x+1}{x-1}\right)+2(x-1+2(x+1))=-2(x-1)+\frac{x-1}{x+1}+x^2.\]After some canceling, distributing, and rearranging, we obtain
 \[4\left(\frac{x+1}{x-1}\right)=x^2-8x-4.\]Multiplying both sides by $x-1$ results in
 \[4x+4=x^3-x^2-8x^2+8x-4x+4,\]which can be rearranged into
 \[x^3-9x^2=0\]and factored into
 \[x^2(x-9)=0.\]This means that $x$ equals $0$ or $9$, and since a side length of $0$ cannot exist, $x=9$.

 
As a result, the triangle must have sides in the ratio of $9:2:9$. Since the triangle must have integer side lengths, and these values share no common factors greater than $1$, the triangle with the smallest possible perimeter under these restrictions has a perimeter of $9+2+9=\boxed{020}$. ~\href{https://artofproblemsolving.com/wiki/index.php?title=User:Emerald\_block}{emerald\_block}

 





\textbf{Solution 3 (Various Techniques)}

Before we start thinking about the problem, let’s draw it out;

 
                    \begin{figure}[H]
                        \centering
                        \includegraphics[width=0.5\textwidth]{Solution/AIME_2019_I_11_Solution_Figure_2.png}
                    \end{figure}
                    

 For the sake of space, I've drawn only 2 of the 3 excircles because the third one looks the same as the second large one because the triangle is isosceles. By the incenter-excenter lemma, $AII_A$ and $BII_B$ are collinear, $E$ is the tangent of circle $I_B$ to $AC$, $F$ is the tangent of that circle to the extension of $BC$, and $J$ is the tangent of the circle to the extension of $BA$. The interesting part of the diagram is circle $\omega$, which is internally tangent to circle $I_A$ yet externally tangent to circle $I_B$. Therefore, perhaps we can relate the radius of this circle to the semiperimeter of triangle $ABC$. 

 We can see that the radius of circle $\omega$ is $2r_{I_A}+r$ using the incenter and A-excenter of our main triangle. This radius is also equal to $BI_B - BI - r_{I_B}$ from the incenter and B-excenter of our triangle. Thus, we can solve for each of these separately in terms of the lengths of the triangle and set them equal to each other to form an equation.

 To find the left hand side of the equation, we have to first find $r$ and $r_{I_A}$. Let $a = AB = AC, b = BD = DC,$ and $h = AD$. Then since the perimeter of the triangle is $2a+2b$, the semiperimeter is $a+b$.

 Now let's take a look at triangle $BDI$. Because $BI$ is the angle bisector of $\angle B$, by the angle bisector theorem, $\frac{AI}{ID} = \frac{BA}{BD} \implies \frac{h-r}{r} = \frac{a}{b}$. Rearranging, we get $r = \frac{hb}{a+b}$.

 Take a look at triangle $AGI$. $AG = a - GB = a - BD = a-b$, $AI = h-r = \frac{ha}{a+b}$ (angle bisector theorem), and $GI = r = \frac{hb}{a+b}$. Now let's analyze triangle $AHI_A$. $AH = AB + BH = AB+ BD = a+b$, $AI_A = h+r_{I_A}$, and $HI_A = r_{I_A}$. Since $\angle GAI = \angle HAI_A$ and $\angle IGA = \angle I_AHA = 90^{\circ}$, triangle $AGI$ and $AHI_A$ are similar by AA. Then $\frac{r_{I_A}}{r} = \frac{h+r_{I_A}}{h-r} \implies r_{I_A} = r \cdot \frac{h+r_{I_A}}{h-r} = \frac{hb}{a+b} \cdot \frac{h+r_{I_A}}{\frac{ha}{a+b}} = \frac{b(h+r_{I_A})}{a}$. Now, solving yields $r_{I_A} = \frac{hb}{a-b}$.

 Finally, the left hand side of our equation is 
 \[\frac{2hb}{a-b} + \frac{hb}{a+b}\]

 Now let's look at triangle $BFI_B$. How will we find $BI_B$? Let's first try to find $BF$ and $I_BF$ in terms of the lengths of the triangle. We recognize:

 $BF = BC + CF = BC + CK$. We really want to have $CA$ instead of $CK$, and $AK$ looks very similar in length to $DC$, so let's try to prove that they are equal. 

 $BJ = BF$, so we can try to add these two and see if we get anything interesting. We have:$BJ + BF = BA + AJ + BC + CF = BA + AE + BC + CE = BA + BC + CA$, which is our perimeter. Thus, $BF = a+b$.

 Triangle $BDI$ is similar to triangle $BFI_B$ by AA, and we know that $BD = b$, and $ID = r = \frac{hb}{a+b}$, so thus $I_BF = BF \cdot \frac{ID}{BD} = (a+b) \cdot \frac{\frac{hb}{a+b}}{b} = \frac{hb}{b} = h$. Thus, the height of this triangle is $h$ by similarity ratios, the same height as vertex $A$. By the Pythagorean Theorem, $BI_B = \sqrt{(a+b)^2 + h^2}$ and by similarity ratios, $BI = \frac{b}{a+b} \cdot \sqrt{(a+b)^2 + h^2}$. Finally, $r_{I_B} = I_BF = h$, and thus the right hand side of our equation is 
 \[\sqrt{(a+b)^2 + h^2} - \frac{b}{a+b} \cdot \sqrt{(a+b)^2 + h^2} - h = \sqrt{(a+b)^2 + h^2}(1 - \frac{b}{a+b}) - h = \sqrt{(a+b)^2 + h^2} \cdot \frac{a}{a+b} - h\].

 Setting the two equal, we have 
 \[\frac{2hb}{a-b} + \frac{hb}{a+b} = \sqrt{(a+b)^2 + h^2} \cdot \frac{a}{a+b} - h\]

 Multiplying both sides by $(a+b)(a-b)$ we have $2hb(a+b) + hb(a-b) = \sqrt{(a+b)^2 + h^2} \cdot a(a-b) - h(a^2 - b^2)$

 From here, let $b = 1$ arbitrarily; note that we can always scale this value to fit the requirements later. Thus our equation is $2h(a+1) + h(a-1) = \sqrt{(a+1)^2 + h^2} \cdot a(a-1) - h(a^2 - 1)$. Now since $h = \sqrt{a^2 - b^2}$, we can plug into our equation:

 $2h(a+1) + h(a-1) = \sqrt{a² + 2a + 1 + a^2 - b^2 } \cdot a(a-1) - h(a^2 - 1)$. Remembering $b = 1$;

 $\implies 2h(a+1) + h(a-1) = \sqrt{2a^2 + 2a} \cdot a(a-1) - h(a^2 - 1)$

 $\implies 3ha + h + ha^2 - h = \sqrt{2a^2 + 2a} \cdot a(a-1)$

 $\implies ah(3+a) = a(a-1) \cdot \sqrt{2a^2 + 2a}$

 $\implies h^2(3+a)^2 = (a-1)^2 \cdot 2a(a+1)$

 $\implies (a^2 - 1) (3 + a)^2 = 2a(a+1)(a - 1)^2$

 $\implies (3+ a)^2 = 2a(a-1)$

 $\implies 9 + 6a + a^2 = 2a^2 - 2a$

 $\implies a^2 - 8a - 9 = 0$

 $\implies (a-9)(a+1) = 0$

 $\implies a = 9$ because the side lengths have to be positive numbers. Furthermore, because our values for $a$ and $b$ are relatively prime, we don't have to scale down our triangle further, and we are done. Therefore, our answer is $2a + 2b = 18 + 2 = \boxed{020}$

 





\textbf{Solution 4 (Not that hard construction)}

Notice that the $A$-excircle would have to be very small to fit the property that it is internally tangent to $\omega$ and the other two excircles are both externally tangent, given that circle $\omega$'s centre is at the incenter of $\triangle ABC$. If $BC=2$, we see that $AB=AC$ must be somewhere in the $6$ to $13$ range. If we test $6$ by construction, we notice the $A$-excircle is too big for it to be internally tangent to $\omega$ while the other two are externally tangent. This means we should test $8$ or $9$ next. I actually did this and found that $9$ worked, so the answer is $2+9+9=\boxed{20}$. Note that $BC$ cannot be $1$ because then $AB=AC$ would have to be $4.5$ which is not an integer. 

 





\textbf{Solution 5 (Standard geometry)}


                    \begin{figure}[H]
                        \centering
                        \includegraphics[width=0.5\textwidth]{Solution/AIME_2019_I_11_Solution_Figure_3.png}
                    \end{figure}
                    Let $M$ be the midpoint $BC, BM = a, AB= BC = b,$$s = b+a$ be the semiperimeter, $r$ be the inradius. Let $I_A, I_B$ be excenters,  $r_A, r_B$ be exradius, $R$ be radius $\omega.$Then $R = r + 2 r_A,$
 \[r = \sqrt{\frac{(s-2a)(s-b)^2}{s}} = a \sqrt{\frac{b-a}{b+a}},\]
 \[r \cdot s = r_A \cdot (s-2a) \implies r_A =  a \sqrt{\frac{b+a}{b–a}},\]
 \[r \cdot s = r_B \cdot (s-b) \implies r_B =  \sqrt{b^2 – a^2} = AM.\]
 \[II_B =R + r_B = r + 2r_A + r_B \implies\]
 \[II_B=   a \sqrt{\frac{b-a}{b+ a}}+ 2a  \sqrt{\frac{b+a}{b – a}} + \sqrt{b^2-a^2} = b\frac{3a+b}{\sqrt{b^2-a^2}}.\]
 \[\overline{AI_B}= \frac{\overline{A} \cdot 2a - \overline{B} \cdot b + \overline{C}\cdot b}{2a -b + b}-\overline{A} =b \frac{\overline{C}- \overline{B}}{2a},\]
 \[AI_B = b, \overline{AI_B}\perp AI \implies II_B = \sqrt{AI_B^2 + (AM – r)^2},\]
 \[II_B =  \sqrt{b^2 +\left(\sqrt{b^2-a^2}-  a \sqrt{\frac{b-a}{b+ a}}\right)^2} = b\sqrt{\frac{2b}{b+a}}\] Therefore we get problem’s condition in the form of
 \[b\frac{b +3a}{\sqrt{b^2-a^2}} = b \sqrt{\frac {2b}{b+a}} \implies  b + 3a = \sqrt{2b(b-a)} \implies (b-9a)(b+a) = 0 \implies b = 9a.\]

 We use $a = 1$ an get $b = 9, 2s = 18+2 = \boxed{020}$.

 \textbf{vladimir.shelomovskii@gmail.com, vvsss}

 





\textbf{Solution 6 (Trig)}

Let $M$ be the midpoint of $BC$, $I$ be the incenter, $MC = a$, $\angle MCI = \theta$,$r = a\tan{\theta}$ be the inradius, $I_A$, $I_B$ be excenters, and $r_A, r_B$ be exradii.

 By the Incenter-Excenter Lemma, the angles $\angle ICI_A$, $\angle ICI_B$, and $\angle IAI_B$ are all right angles.We can then find that $r_A = a\cot{\theta}$ and $r_B = MA = a\tan{2\theta}$.We also know $\angle CII_B = 2\theta$ since it is an external angle of $\triangle BCI$.Then, $II_B = IC\sec{2\theta} = a\sec{\theta}\sec{2\theta}$. From the condition in the problem, we know that $II_B = r+2r_A+r_B$, so plugging in terms and simplifying gives

 
 \[a\sec{\theta}\sec{2\theta} = a\tan{\theta}+2a\cot{\theta}+a\tan{2\theta}\]
 \[(\tan{\theta}+2\cot{\theta}+\tan{2\theta})(\cos{\theta})(\cos{2\theta})=1\]
 \[(2\cot{\theta}-\tan{\theta})(\cos{\theta})=1\]
 \[\frac{2\cos^2{\theta}}{\sin{\theta}}-\sin{\theta}=1\]
 \[2\cos^2{\theta}-\sin^2{\theta}=\sin{\theta}\]
 \[2-3\sin^2{\theta}=\sin{\theta}\]
 \[(3\sin{\theta}-2)(\sin{\theta}+1)=0\]

 which gives $\sin{\theta}=\frac{2}{3}$. Then $AC = \frac{a}{\cos{2\theta}} = \frac{a}{1-2(\frac{2}{3})^2} = 9a$.Thus, the smallest value is when $a=1$, which gives $\boxed{020}$.

 





\textbf{Solution 7 (Similar to Solution 3)}

Let the incenter of $\triangle ABC$ be $I$ and the centers of the excircles tangent to $BC$ and $AC$ be $O_1$, $O_2$, respectively. Denote the lengths of $AB$ and $AC$ as $x$, the length of $BC$ as $y$, the area of $\triangle ABC$ as $S$, and the semi-perimeter of which as $p$. Let $H_1$ be the point where the incircle of $\triangle ABC$ is tangent to side $BC$, and $H_2$ be the point where the excircle with center $O_2$ is tangent to the extension of side $BC$.

 From the definition of perimeter, it is obvious that $p=\frac{2x+y}{2}$. Apply Heron's formula to yield $S=\sqrt{p(p-x)(p-x)(p-y)}$. Since $I$ is the incenter of $\triangle ABC$, $IH_1$, the radius of the incircle of $\triangle ABC$ is easily $\frac{S}{p}$. Since the circle with center $O_2$ is an excircle of the triangle, $O_2H_2$, the radius of the excircle is easily $\frac{S}{p-x}$. It is also simple to find out that the radius of circle $\omega$ (because of tangency) is equal to the sum of diameter of the excircle with center $O_1$ and the radius of the incircle. Denote the radius of circle $\omega$ as $r$, we can see that: $r=\frac{S}{p}+\frac{2S}{p-y}$. Thus $IO_2=\frac{S}{p}+\frac{2S}{p-y}+\frac{S}{p-x}$.

 Since $IH_1$ and $O_2H_2$ are both tangent to line $BC$, we can draw from similar triangles the conclusion that:

 
 \[\frac{O_2H_2}{IH_1}=\frac{BO_2}{BI}=\frac{IO_2+BI}{BI}=1+\frac{IO_2}{BI}\]

 Due to the properties of isosceles triangles, it is true that $BH_1=\frac{BC}{2}=\frac{y}{2}$. Since $IH_1=\frac{S}{p}$, from the Pythagorean theorem, we can obtain that $BI=\sqrt{\left(\frac{y}{2}\right)^2+\left(\frac{S}{p}\right)^2}$. After plugging in $S$ and $p$ in terms of $x$ and $y$ and further algebraic procedures, the equation turns into:

 
 \[4x^2-16xy-9y^2=0\]

 Solve $x$ in terms of $y$ and eliminate the negative root to get:

 
 \[x=4.5y\]

 Therefore, the minimum of the perimeter of $\triangle ABC$ with integer side lengths is $2\times9+2=\boxed{020}$

 


 \href{https://mathproblemsolvingskills.wordpress.com/}{MathProblemSolvingSkills.com}

 


 This solution is the video solution for Solution 3 - not posted by ~KingRavi

 




	\item (\textit{AIME I 2019}) Given $f(z) = z^2-19z$, there are complex numbers $z$ with the property that $z$, $f(z)$, and $f(f(z))$ are the vertices of a right triangle in the complex plane with a right angle at $f(z)$. There are positive integers $m$ and $n$ such that one such value of $z$ is $m+\sqrt{n}+11i$. Find $m+n$.


 
 \textbf{Answer: }230



\textbf{Solution 1}

Notice that we must have 
 \[\frac{f(f(z))-f(z)}{f(z)-z}=-\frac{f(f(z))-f(z)}{z-f(z)}\in i\mathbb R .\]However, $f(t)-t=t(t-20)$, so
 \begin{align*} \frac{f(f(z))-f(z)}{f(z)-z}&=\frac{(z^2-19z)(z^2-19z-20)}{z(z-20)}\\ &=\frac{z(z-19)(z-20)(z+1)}{z(z-20)}\\ &=(z-19)(z+1)\\ &=(z-9)^2-100. \end{align*}
Then, the real part of $(z-9)^2$ is $100$. Since $\text{Im}(z-9)=\text{Im}(z)=11$, let $z-9=a+11i$. Then, 
 \[100=\text{Re}((a+11i)^2)=a^2-121\implies a=\pm\sqrt{221}.\]It follows that $z=9+\sqrt{221}+11i$, and the requested sum is $9+221=\boxed{230}$.

 (Solution by TheUltimate123)

 





\textbf{Solution 2}

We will use the fact that segments $AB$ and $BC$ are perpendicular in the complex plane if and only if $\frac{a-b}{b-c}\in i\mathbb{R}$. To prove this, note that when dividing two complex numbers you subtract the angle of one from the other. Therefore, if the two complex numbers are perpendicular, the difference between their arguments will be 90 degrees, so subtracting the angles will yield an imaginary number with no real part(an argument of 90 degrees puts a complex number on the imaginary axis).

 Now to apply this: 
 \[\frac{f(z)-z}{f(f(z))-f(z)}\in i\mathbb{R}\]
 \[\frac{z^2-19z-z}{(z^2-19z)^2-19(z^2-19z)-(z^2-19z)}\]
 \[\frac{z^2-20z}{z^4-38z^3+341z^2+380z}\]
 \[\frac{z(z-20)}{z(z+1)(z-19)(z-20)}\]
 \[\frac{1}{(z+1)(z-19)}\in i\mathbb{R}\]

 The factorization of the nasty denominator above is made easier with the intuition that $(z-20)$ must be a divisor for the problem to lead anywhere. Now we know $(z+1)(z-19)\in i\mathbb{R}$ so using the fact that the imaginary part of $z$ is $11i$ and calling the real part r, 

 
 \[(r+1+11i)(r-19+11i)\in i\mathbb{R}\]
 \[r^2-18r-140=0\]

 solving the above quadratic yields $r=9+\sqrt{221}$ so our answer is $9+221=\boxed{230}$

 





\textbf{Solution 3}

I would like to use a famous method, namely the coni method. 

 \begin{verbatim}
Statement .If we consider three complex numbers $A,B,C$ in argand plane then $\angle ABC =\arg{\frac{B-A}{B-C}}$.
\end{verbatim}According to the question given, we can assume ,$A= f(f(z)),B=f(z),C= z$ respectively.

 WLOG,$Z_1= \frac{f(f(z))-f(z)}{z-f(z)}$.According to the question $\arg{Z_1}=\frac{\pi}{2}$.

 So,$\Re (Z_1)=0$.

 Now, $Z_1=\frac{z(z-19)(z+1)(z-20)}{-z(z-20)}$.

 $\implies Z_1= -(z^2-18z-19)$.WLOG,$z=a+11i$ .where $a=m+\sqrt{n}$.

 So,$\Re (Z_1)= -(a^2-18a-140)$.Solving,$a^2-18a-140=0$ .get ,

 $a=9$±$\sqrt{221}$.So, possible value of $a=9+\sqrt{221}$.

 $m+n=\boxed{230}$.~ftheftics.

 \textbf{Proof of this method} Note that if we translate a triangle, the measures of all of its angles stay the same. So we can translate $ABC$ on the complex plane so that $B=0$. Let the images of $A, B, C$ be $A', B', C'$ respectively. Then, we can use the formula:

 
 \[re^{i\theta}=r\cos(\theta)+i \cdot r\sin(\theta)\] (This is known as Euler's Formula.)

 Using Euler's formula (represent each complex number in polar form, then use exponent identities), we can show that 
 \[\text{arg}(\frac{A'}{C'})=\text{arg}(A')-\text{arg}(C')=\angle(A'B'C')=\angle(ABC)\] So this method is valid.~Math4Life2020

 





\textbf{Solution 4}

It is well known that $AB$ is perpendicular to $CD$ iff $\frac{d-c}{b-a}$ is a pure imaginary number. Here, we have that $A=z$, $B,C=f(z)$, and $D=f(f(z))$. This means that this is equivalent to $\frac{f(f(z))-f(z)}{f(z)-z}$ being a pure imaginary number. Plugging in $f(z)=z^2-19z$, we have that $\frac{(z^2-19z)-19(z^2-19z)-(z^2-19z)}{z^2-19z-z}$ being pure imaginary. Factoring and simplifying, we find that this is simply equivalent to $(z-19)(z+1)$ being pure imaginary. We let $z=a+bi$, so this is equivalent to $(a+bi-19)(a+bi+1)$ being pure imaginary. Expanding the product, this is equivalent to $a^2+abi+a+abi-b^2+bi-19a-19bi-19$ being pure imaginary. Taking the real part of this, and setting this equal to $0$, we have that $a^2-18a-b^2-19=0$. Since $b=11$, we have that $a^2-18a-140=0$. By the quadratic formula, $a=9 \pm \sqrt{221}$, and taking the positive root gives that $a=9+ \sqrt{221}$, so the answer is $9+221=230$

 





\textbf{Solution 5 (Official MAA)}

The arguments of the two complex numbers differ by $90^\circ$ if the ratio of the numbers is a pure imaginary number. Thus three distinct complex numbers $A,\,B,$ and $C$ form a right triangle at $B$ if and only if $\frac{C-B}{B-A}$ has real part equal to $0.$ Hence 
 \begin{align*} \frac{f(f(z))-f(z)}{f(z)-z}&=\frac{(z^2-19z)^2-19(z^2-19z)-(z^2-19z)}{(z^2-19z)-z}\\ &=\frac{(z^2-19z)(z^2-19z-19-1)}{z^2-20z}\\ &=\frac{z(z-19)(z+1)(z-20)}{z(z-20)} \\ &=z^2-18z-19 \end{align*}
 must have real part equal to $0.$ If $z=x+11i,$ the real part of $z^2-18z-19$ is $x^2-11^2-18x-19,$ which is $0$ when $x=9\pm\sqrt{221}.$ The requested sum is $9+221=\boxed{230}.$

 





\textbf{Solution 6 (Dot product)}

Firstly, the angle between the three complex numbers is equivalent to the angle between $f(z)-z,0,$ and $f(f(z))-f(z)$.

 Using $f(z)=z(z-19)$ to help expand,

 $z-f(z)=20z-z^2$ and $f(f(z))-f(z)=(z^2-19z)(z^2-19z-19)-(z^2-19z)$.

 The second equation can be rewritten as $(z^2-19z)(z^2-19z-20)=z(z-19)(z-20)(z+1)$

 Note that the angle between $(a+bi), 0,$ and $(a+bi)(c+di)$ is the same as the angle between $1, 0,$and $c+di$.The proof of this is as follows:

 Treating the complex numbers like vectors (e.g. $a+bi$ turns into $\left[\begin{array}{c}  a \\ b \end{array}\right]$) we have $cos \theta = \frac{a \cdot (ac-bd) + b \cdot (ad+bc)}{2\sqrt{a^2+b^2}\sqrt{a^2+b^2}\sqrt{c^2+d^2}}$

 $cos \theta = \frac{a^2 c - abd + abd + b^2 c}{2(a^2+b^2)\sqrt{c^2+d^2}}$

 $cos \theta = \frac{(a^2+b^2)c}{2(a^2+b^2)\sqrt{c^2+d^2}}=\frac{c}{2\sqrt{c^2+d^2}}$.

 Using the dot product formula for the cosine of the angle between the other two vectors (1, c + di) we get the same result. Thus, it is proven.

 Now, we see that the three resulting complex numbers have the same angle as $-1,0,$ and $(z-19)(z+1)$ as a factor of $z(z-20)$ can be taken out of both expressions. Note that the value of $z$ clearly is not 0, nor is it 20 as it has a value of $11i$.

 Expanding $z^2-18z-19$ with $z=a+11i$ yields

 $a^2-121+22ai-18a-198i-19=a^2-18a-140+22ai-198i$. Turning these into vectors again, the resulting relation is

 $cos(90^\circ)=0=\frac{-(a^2-18a-140)}{2\sqrt{(a^2-18a-140)^2 + (22a-198)^2} \cdot 1}$

 In order for this thing to be 0, the numerator must be 0, so setting it equal to 0 yields

 $a^2-18a-140=0 \Rightarrow a=9\pm \sqrt{81+140}$. Neither of these make the denominator in the $cos(90^\circ)$ expression 0, so they are valid. The requested result is the positive root, so the answer is $\boxed{230}$.

 




	\item (\textit{AIME I 2019}) Triangle $ABC$ has side lengths $AB=4$, $BC=5$, and $CA=6$. Points $D$ and $E$ are on ray $AB$ with $AB<AD<AE$. The point $F \neq C$ is a point of intersection of the circumcircles of $\triangle ACD$ and $\triangle EBC$ satisfying $DF=2$ and $EF=7$. Then $BE$ can be expressed as $\frac{a+b\sqrt{c}}{d}$, where $a$, $b$, $c$, and $d$ are positive integers such that $a$ and $d$ are relatively prime, and $c$ is not divisible by the square of any prime. Find $a+b+c+d$.


 
 \textbf{Answer: }032



\textbf{Solution 1}


                    \begin{figure}[H]
                        \centering
                        \includegraphics[width=0.5\textwidth]{Solution/AIME_2019_I_13_Solution_Figure_1.png}
                    \end{figure}
                    

 Notice that 
 \[\angle DFE=\angle CFE-\angle CFD=\angle CBE-\angle CAD=180-B-A=C.\]By the Law of Cosines, 
 \[\cos C=\frac{AC^2+BC^2-AB^2}{2\cdot AC\cdot BC}=\frac34.\]Then, 
 \[DE^2=DF^2+EF^2-2\cdot DF\cdot EF\cos C=32\implies DE=4\sqrt2.\]Let $X=\overline{AB}\cap\overline{CF}$, $a=XB$, and $b=XD$. Then, 
 \[XA\cdot XD=XC\cdot XF=XB\cdot XE\implies b(a+4)=a(b+4\sqrt2)\implies b=a\sqrt2.\]However, since $\triangle XFD\sim\triangle XAC$, $XF=\frac{4+a}3$, but since $\triangle XFE\sim\triangle XBC$, 
 \[\frac75=\frac{4+a}{3a}\implies a=\frac54\implies BE=a+a\sqrt2+4\sqrt2=\frac{5+21\sqrt2}4,\]and the requested sum is $5+21+2+4=\boxed{032}$.

 (Solution by TheUltimate123)

 





\textbf{Solution 2}

Define $\omega_1$ to be the circumcircle of $\triangle ACD$ and $\omega_2$ to be the circumcircle of $\triangle EBC$.

 Because of exterior angles, 

 $\angle ACB = \angle CBE - \angle CAD$

 But $\angle CBE = \angle CFE$ because $CBFE$ is cyclic. In addition, $\angle CAD = \angle CFD$ because $CAFD$ is cyclic. Therefore, $\angle ACB = \angle CFE - \angle CFD$. But $\angle CFE - \angle CFD = \angle DFE$, so $\angle ACB = \angle DFE$. Using Law of Cosines on $\triangle ABC$, we can figure out that $\cos(\angle ACB) = \frac{3}{4}$. Since $\angle ACB = \angle DFE$, $\cos(\angle DFE) = \frac{3}{4}$. We are given that $DF = 2$ and $FE = 7$, so we can use Law of Cosines on $\triangle DEF$ to find that $DE = 4\sqrt{2}$.

 Let $G$ be the intersection of segment $\overline{AE}$ and $\overline{CF}$. Using Power of a Point with respect to $G$ within $\omega_1$, we find that $AG \cdot GD = CG \cdot GF$. We can also apply Power of a Point with respect to $G$ within $\omega_2$ to find that $CG \cdot GF = BG \cdot GE$. Therefore, $AG \cdot GD = BG \cdot GE$.

 
 \[AG \cdot GD = BG \cdot GE\]
 \[(AB + BG) \cdot GD = BG \cdot (GD + DE)\]
 \[AB \cdot GD + BG \cdot GD = BG \cdot GD + BG \cdot DE\]
 \[AB \cdot GD = BG \cdot DE\]
 \[4 \cdot GD = BG \cdot 4\sqrt{2}\]
 \[GD = BG \cdot \sqrt{2}\]

 Note that $\triangle GAC$ is similar to $\triangle GFD$. $GF = \frac{BG + 4}{3}$. Also note that $\triangle GBC$ is similar to $\triangle GFE$, which gives us $GF = \frac{7 \cdot BG}{5}$. Solving this system of linear equations, we get $BG = \frac{5}{4}$. Now, we can solve for $BE$, which is equal to $BG(\sqrt{2} + 1) + 4\sqrt{2}$. This simplifies to $\frac{5 + 21\sqrt{2}}{4}$, which means our answer is $\boxed{032}$.

 





\textbf{Solution 3}

Construct $FC$ and let $FC\cap AE=K$. Let $FK=x$. Using $\triangle FKE\sim \triangle BKC$, 
 \[BK=\frac{5}{7}x\] Using $\triangle FDK\sim ACK$, it can be found that 
 \[3x=AK=4+\frac{5}{7}x\to x=\frac{7}{4}\] This also means that $BK=\frac{21}{4}-4=\frac{5}{4}$. It suffices to find $KE$. It is easy to see the following: 
 \[180-\angle ABC=\angle KBC=\angle KFE\] Using reverse Law of Cosines on $\triangle ABC$, $\cos{\angle ABC}=\frac{1}{8}\to \cos{180-\angle ABC}=\frac{-1}{8}$. Using Law of Cosines on $\triangle EFK$ gives $KE=\frac{21\sqrt 2}{4}$, so $BE=\frac{5+21\sqrt 2}{4}\to \boxed{\textbf{032}}$.-franchester

 





\textbf{Solution 4 (No <C = <DFE, no LoC)}

Let <img alt="$P=AE\cap CF$" class="latex" height="13" src="//latex.artofproblemsolving.com/5/e/0/5e08bc35f5474607458539d2cfb49c2e558030c8.png" style="vertical-align: 0px" width="115"/>. Let $CP=5x$ and $BP=5y$; from $\triangle{CBP}\sim\triangle{EFP}$ we have $EP=7x$ and $FP=7y$. From $\triangle{CAP}\sim\triangle{DFP}$ we have $\frac{6}{4+5y}=\frac{2}{7y}$ giving $y=\frac{1}{4}$. So $BP=\frac{5}{4}$ and $FP=\frac{7}{4}$. These similar triangles also gives us $DP=\frac{5}{3}x$ so $DE=\frac{16}{3}x$. Now, Stewart's Theorem on $\triangle{FEP}$ and cevian $FD$ tells us that 
 \[\frac{560}{9}x^3+28x=\frac{49}{3}x+\frac{245}{3}x,\]so $x=\frac{3\sqrt{2}}{4}$. Then $BE=\frac{5}{4}+7x=\frac{5+21\sqrt{2}}{4}$ so the answer is $\boxed{032}$ as desired. (Solution by Trumpeter, but not added to the Wiki by Trumpeter)

 





\textbf{Solution 5}

Connect $CF$ meeting $AE$ at $J$. We can observe that $\triangle{ACJ}\sim \triangle{FJD}$ Getting that $\frac{AJ}{FJ}=\frac{AC}{FD}=3$. We can also observe that $\triangle{CBJ}\sim \triangle{EFJ}$, getting that $\frac{CB}{EF}=\frac{BJ}{FJ}=\frac{CJ}{EJ}=\frac{5}{7}$

 Assume that $BJ=5x;FJ=7x$, since $\frac{AJ}{FJ}=3$, we can get that $\frac{AJ}{FJ}=\frac{AB+BJ}{FJ}=\frac{4+5x}{7x}=3$, getting that $x=\frac{1}{4}; BJ=\frac{5}{4}; FJ=\frac{7}{4}$

 Using Power of Point, we can get that $BJ * EJ=CJ*FJ; DJ * AJ=CJ * FJ$ Assume that $DJ=5k,CJ=15k$, getting that $JE=21k, DE=16k$

 Now applying Law of Cosine on two triangles, $\triangle{ACJ};\triangle{FJE}$ separately, we can get two equations

 $(1): (15k)^2+(\frac{21}{4})^2-2*15k *\frac{21}{4} * cos\angle{CJA}=36$

 $(2):(21k)^2+(\frac{7}{4})^2-2*\frac{7}{4} * 21k*cos\angle{FJE}=49$

 Since $\angle{CJA}=\angle{FJE}$, we can use $15(2)-7(1)$ to eliminate the $cos$ term

 Then we can get that $5040k^2=630$, getting $k=\frac{\sqrt{2}}{4}$

 $BE=21k=\frac{21\sqrt{2}}{4}; BJ=\frac{5}{4}$, so the desired answer is $\frac{21\sqrt{2}+5}{4}$, which leads to the answer $\boxed{032}$

 





\textbf{Solution 6}

First, let $AE$ and $CF$ intersect at $X$. Our motivation here is to introduce cyclic quadrilaterals and find useful relationships in terms of angles. Observe that
 \[\angle DFE = \angle XFE - \angle XFD = \angle CBE - \angle CAB = 180 - \angle ABC - \angle CAB = \angle BAC\]By the so-called "Reverse Law of Cosines" on $\triangle ABC$ we have
 \[\cos(\angle BAC) = \frac{4^2 - 5^2 - 6^2}{-2 \cdot 5 \cdot 6} = \frac{3}{4}\]Applying on $\triangle DFE$ gives
 \[DE^2 = 2^2 + 7^2 - 2 \cdot 2 \cdot 7 \cos(\angle DFE)\]
 \[= 2^2 + 7^2 - 2 \cdot 2 \cdot 7 \cdot \frac{3}{4}\]
 \[=32\]So $DE = 4 \sqrt{2}$, now by our cyclic quadrilaterals again, we are motivated by the multiple appearances of similar triangles throughout the figure. We want some that are related to $BX$ and $XD$, which are crucial lengths in the problem. Suppose $BX = r, XD = s$ for simplicity. We have:

 
 \[\triangle AXC \sim \triangle FXD\]and
 \[\triangle BXC \sim \triangle FXE\]

 So
 \[\frac{AX}{FX} = \frac{XC}{XD} = \frac{AC}{FD} \implies \frac{4 + r}{FX} = \frac{XC}{s} = 3\]
 \[\frac{BX}{FX} = \frac{XC}{XE} = \frac{BC}{FE} \implies \frac{r}{FX} = \frac{XC}{s + 4 \sqrt{2}} = \frac{5}{7}\]
 \[\implies \frac{4 + r}{r} = \frac{s + 4 \sqrt{2}}{s} = \frac{21}{5}\]
 \[\implies r = \frac{5}{4}, s = \frac{5 \sqrt{2}}{4}\]So $BE = r + s + 4 \sqrt{2} = \frac{5 + 21 \sqrt{2}}{4}$. The requested sum is $5 + 21 + 2 + 4 = \boxed{032}$.

 




	\item (\textit{AIME I 2019}) Find the least odd prime factor of $2019^8 + 1$.


 
 \textbf{Answer: }097



\textbf{Solution}

We know that $2019^8 \equiv -1 \pmod{p}$ for some prime $p$. We want to find the smallest odd possible value of $p$. By squaring both sides of the congruence, we find $2019^{16} \equiv 1 \pmod{p}$.   

 Since $2019^{16} \equiv 1 \pmod{p}$, the order of $2019$ modulo $p$ is a positive divisor of $16$.

 However, if the order of $2019$ modulo $p$ is $1, 2, 4,$ or $8,$ then $2019^8$ will be equivalent to $1 \pmod{p},$ which contradicts the given requirement that $2019^8\equiv -1\pmod{p}$.

 Therefore, the order of $2019$ modulo $p$ is $16$. Because all orders modulo $p$ divide $\phi(p)$, we see that $\phi(p)$ is a multiple of $16$. As $p$ is prime, $\phi(p) = p\left(1 - \dfrac{1}{p}\right) = p - 1$. Therefore, $p\equiv 1 \pmod{16}$. The two smallest primes equivalent to $1 \pmod{16}$ are $17$ and $97$. Because $16 | p - 1$, and $p - 1  \geq 16$, each possible value of $p$ must be verified by manual calculation to make sure that $p | 2019^8+1$. As $2019^8 \not\equiv -1 \pmod{17}$ and $2019^8 \equiv -1 \pmod{97}$, the smallest possible $p$ is thus $\boxed{097}$.

 \textbf{Note to solution} $\phi(k)$ is the Euler Totient Function of integer $k$. $\phi(k)$ is the number of positive integers less than $k$ relatively prime to $k$. Define the numbers $k_1,k_2,k_3,\cdots,k_n$ to be the prime factors of $k$. Then, we have
 \[\phi(k)=k\cdot \prod^n_{i=1}\left(1-\dfrac{1}{k_i}\right).\]A property of the Totient function is that, for any prime $p$, $\phi(p)=p-1$.

 Euler's Totient Theorem states that
 \[a^{\phi(k)} \equiv 1\pmod k\]if $\gcd(a,k)=1$.

 Furthermore, the order $a$ modulo $n$ for an integer $a$ relatively prime to $n$ is defined as the smallest positive integer $d$ such that $a^{d} \equiv 1\pmod n$. An important property of the order $d$ is that $d|\phi(n)$. It's important to remember that the order $d \leq \phi(n)$, and Euler's Totient Theorem does not guarantee equality. That's why in this problem we had to test primes such that $d=16 | (p-1)$, and find the smallest such prime $p$.~First

 





\textbf{Solution 2 (Basic Modular Arithmetic)}

In this solution, $k$ will represent an arbitrary nonnegative integer.

 We will show that any potential prime $p$ must be of the form $16k+1$ through a proof by contradiction. Suppose that there exists some prime $p$ that can not be expressed in the form $16k+1$ that is a divisor of $2019^8+1$. First, note that if the prime $p$ is a divisor of $2019$, then $2019^8$ is a multiple of $p$ and $2019^8+1$ is not. Thus, $p$ is not a divisor of $2019$.

 Because $2019^8+1$ is a multiple of $p$, $2019^8+1\equiv0\pmod{p}$. This means that $2019^8\equiv-1\pmod{p}$, and by raising both sides to an arbitrary odd positive integer, we have that $2019^{16k+8}\equiv-1\pmod{p}$.

 Then, since the problem requires an odd prime, $p$ can be expressed as $16k+m$, where $m$ is an odd integer ranging from $3$ to $15$, inclusive. By Fermat's Little Theorem, $2019^{p-1}\equiv1\pmod{p}$, and plugging in values, we get $2019^{16k+n}\equiv1\pmod{p}$, where $n=m-1$ and is thus an even integer ranging from $2$ to $14$, inclusive.

 If $n=8$, then $2019^{16k+8}\equiv1\pmod{p}$, which creates a contradiction. If $n$ is not a multiple of $8$ but is a multiple of $4$, squaring both sides of $2019^{16k+n}\equiv1\pmod{p}$ also results in the same contradictory equivalence. For all remaining $n$, raising both sides of $2019^{16k+n}\equiv1\pmod{p}$ to the $4$th power creates the same contradiction. (Note that $32k$ and $64k$ can both be expressed in the form $16k$.)

 Since we have proved that no value of $n$ can work, this means that a prime must be of the form $16k+1$ in order to be a factor of $2019^8+1$. The smallest prime of this form is $17$, and testing it, we get
 \[2019^8+1\equiv13^8+1\equiv169^4+1\equiv(-1)^4+1\equiv1+1\equiv2\pmod{17},\]so it does not work. The next smallest prime of the required form is $97$, and testing it, we get
 \[2019^8+1\equiv(-18)^8+1\equiv324^4+1\equiv33^4+1\equiv1089^2+1\equiv22^2+1\equiv484+1\equiv-1+1\equiv0\pmod{97},\]so it works. Thus, the answer is $\boxed{097}$. ~\href{https://artofproblemsolving.com/wiki/index.php?title=User:Emerald\_block}{emerald\_block}

 





\textbf{Solution 3 (Official MAA)}

Suppose prime $p>2$ divides $2019^8+1.$ Then $2019^8\equiv -1\pmod p.$ Squaring gives $2019^{16}\equiv 1\pmod p.$ If $2019^m\equiv 1 \pmod p$ for some $0<m<16,$ it follows that 
 \[2019^{\gcd(m,16)}\equiv 1\pmod p.\] But $2019^8\equiv -1\pmod p,$ so $\gcd(m,16)$ cannot divide $8,$ which is a contradiction. Thus $2019^{16}$ is the least positive power congruent to $1\pmod p.$ By Fermat's Little Theorem, $2019^{p-1}\equiv 1\pmod p.$ It follows that $p=16k+1$ for some positive integer $k.$ The least two primes of this form are $17$ and $97.$ The least odd factor of $2019^8+1$ is not $17$ because 
 \[2019\equiv 13\pmod{17}\qquad \text{and}\qquad 13^2\equiv 169\equiv -1\pmod{17},\] which implies $2019^8\equiv 1\not\equiv -1\pmod {17}.$ But $2019\equiv -18\pmod{97},$ so 
 \begin{align*} (-18)^2=324&\equiv 33\pmod{97}, \\ 33^2=1089&\equiv 22\pmod{97},\,\text{and} \\ 22^2=484&\equiv -1\pmod{97}.\end{align*}
 Thus the least odd prime factor is $97.$

 In fact, $2019^8+1=2\cdot97\cdot p,$ where $p$ is the $25$-digit prime 
 \[1423275002072658812388593.\]

 


 





\textbf{Solution 4}

Let $n$ represent the least odd prime that the question is asking for. 

 We can see that $2019^8\equiv -1\pmod n.$

 Squaring both sides we get $2019^{16}\equiv 1\pmod n.$

 From Fermat's Little Theorem $a^{n-1}\equiv 1\pmod n$, we know that $n-1$ has to be a multiple of 16. So we want to test prime numbers that fit this. 

 The first prime we get is 17 and when we try it in $2019^8\equiv -1\pmod n,$

 $2019^{8}\equiv 13^8\pmod {17}$

 $169^{4}\equiv 16^4\pmod {17}$

 $256^2\equiv 1\pmod {17}$

 $1+1=2$ 

 We see that 17 does not work. The next number that works from the application of Fermat's Little Theorem is 97 and when we try that,

 $2019^{8}\equiv 79^8\pmod {97}$

 $6241^{4}\equiv 33^4\pmod {97}$

 $1089^2\equiv 22^2\pmod {97}$

 $484\equiv -1\pmod {97}$

 $-1+1=0$

 Thus our answer is $\boxed{097}$. ~pengf

 On The Spot STEM:

 


 




	\item (\textit{AIME I 2019}) Let $\overline{AB}$ be a chord of a circle $\omega$, and let $P$ be a point on the chord $\overline{AB}$. Circle $\omega_1$ passes through $A$ and $P$ and is internally tangent to $\omega$. Circle $\omega_2$ passes through $B$ and $P$ and is internally tangent to $\omega$. Circles $\omega_1$ and $\omega_2$ intersect at points $P$ and $Q$. Line $PQ$ intersects $\omega$ at $X$ and $Y$. Assume that $AP=5$, $PB=3$, $XY=11$, and $PQ^2 = \frac{m}{n}$, where $m$ and $n$ are relatively prime positive integers. Find $m+n$.


 
 \textbf{Answer: }065



\textbf{Solution 1}


                    \begin{figure}[H]
                        \centering
                        \includegraphics[width=0.5\textwidth]{Solution/AIME_2019_I_15_Solution_Figure_1.png}
                    \end{figure}
                    Let $O_1$ and $O_2$ be the centers of $\omega_1$ and $\omega_2$, respectively. There is a homothety at $A$ sending $\omega$ to $\omega_1$ that sends $B$ to $P$ and $O$ to $O_1$, so $\overline{OO_2}\parallel\overline{O_1P}$. Similarly, $\overline{OO_1}\parallel\overline{O_2P}$, so $OO_1PO_2$ is a parallelogram. Moreover, 
 \[\angle O_1QO_2=\angle O_1PO_2=\angle O_1OO_2,\]hence $OO_1O_2Q$ is cyclic. However, 
 \[OO_1=O_2P=O_2Q,\]so $OO_1O_2Q$ is an isosceles trapezoid. Since $\overline{O_1O_2}\perp\overline{XY}$, $\overline{OQ}\perp\overline{XY}$, so $Q$ is the midpoint of $\overline{XY}$.

 By Power of a Point, $PX\cdot PY=PA\cdot PB=15$. Since $PX+PY=XY=11$ and $XQ=11/2$, 
 \[XP=\frac{11-\sqrt{61}}2\implies PQ=XQ-XP=\frac{\sqrt{61}}2\implies PQ^2=\frac{61}4,\]and the requested sum is $61+4=\boxed{065}$.

 (Solution by TheUltimate123)

 \textbf{Note} One may solve for $PX$ first using PoAP, $PX = \frac{11}{2} - \frac{\sqrt{61}}{2}$. Then, notice that $PQ^2$ is rational but $PX^2$ is not, also $PX = \frac{XY}{2} - \frac{\sqrt{61}}{2}$. The most likely explanation for this is that $Q$ is the midpoint of $XY$, so that $XQ = \frac{11}{2}$ and $PQ=\frac{\sqrt{61}}{2}$. Then our answer is $m+n=61+4=\boxed{065}$. One can rigorously prove this using the methods above

 





\textbf{Solution 2 (olympiad harmonic)}

Let the tangents to $\omega$ at $A$ and $B$ intersect at $R$. Then, since $RA^2=RB^2$, $R$ lies on the radical axis of $\omega_1$ and $\omega_2$, which is $\overline{PQ}$. It follows that 
 \[-1=(A,B;X,Y)\stackrel{A}{=}(R,P;X,Y).\]Let $Q'$ denote the midpoint of $\overline{XY}$. By the Midpoint of Harmonic Bundles Lemma(EGMO 9.17), 
 \[RP\cdot RQ'=RX\cdot RY=RA^2=RP\cdot RQ,\]whence $Q=Q'$. Like above, $XP=\frac{11-\sqrt{61}}2$. Since $XQ=\frac{11}2$, we establish that $PQ=\frac{\sqrt{61}}2$, from which $PQ^2=\frac{61}4$, and the requested sum is $61+4=\boxed{065}$.

 (Solution by TheUltimate123)

 





\textbf{Solution 3}

Firstly we need to notice that $Q$ is the middle point of $XY$. Assume the center of circle $w, w_1, w_2$ are $O, O_1, O_2$, respectively. Then $A, O_1, O$ are collinear and $O, O_2, B$ are collinear. Link $O_1P, O_2P, O_1Q, O_2Q$. Notice that, $\angle B=\angle A=\angle APO_1=\angle BPO_2$. As a result, $PO_1\parallel O_2O$ and $OO_1\parallel O_2P$. So we have parallelogram $PO_2OO_1$. So $\angle O_2PO_1=\angle O$ Notice that, $O_1O_2\bot PQ$ and $O_1O_2$ divides $PQ$ into two equal length pieces, So we have $\angle O_2QO_1=\angle O_2PO_1=\angle O$. As a result, $O_2, Q, O, O_1,$ lie on one circle. So $\angle OQO_1=\angle OO_2O_1=\angle O_2O_1P$. Notice that since $\angle O_1PQ+\angle O_2O_1P=90^{\circ}$, we have $\angle OQP=\angle OQO_1 + \angle O_1QP = \angle O_2O_1P + O_1PQ=90^{\circ}$. As a result, $OQ\bot PQ$. So $Q$ is the middle point of $XY$.

 Back to our problem. Assume $XP=x$, $PY=y$ and $x<y$. Then we have $AP\cdot PB=XP\cdot PY$, that is, $xy=15$. Also, $XP+PY=x+y=XY=11$. Solve these above, we have $x=\frac{11-\sqrt{61}}{2}=XP$. As a result, we have $PQ=XQ-XP=\frac{11}{2}-\frac{11-\sqrt{61}}{2}=\frac{\sqrt{61}}{2}$. So, we have $PQ^2=\frac{61}{4}$. As a result, our answer is $m+n=61+4=\boxed{065}$.

 


 Solution By BladeRunnerAUG (Fanyuchen20020715). Edited by bgn4493.

 





\textbf{Solution 4}


                    \begin{figure}[H]
                        \centering
                        \includegraphics[width=0.5\textwidth]{Solution/AIME_2019_I_15_Solution_Figure_2.png}
                    \end{figure}
                    Note that the tangents to the circles at $A$ and $B$ intersect at a point $Z$ on $XY$ by radical axis theorem. Since $\angle ZAB = \angle ZQA$ and $\angle ZBA = \angle ZQB$, we have 
 \[\angle AZB + \angle AQB = \angle AZB + \angle ZAB + \angle ZBA = 180^{\circ},\]so $ZAQB$ is cyclic. 

 But if $O$ is the center of $\omega$, clearly $ZAOB$ is cyclic with diameter $ZO$, so $\angle ZQO = 90^{\circ}$ implies that $Q$ is the midpoint of $XY$. Then, by power of point $P$, 
 \[PY \cdot PX = PA \cdot PB = 15,\] whereas it is given that $PY+PX = 11$. Thus 
 \[PY, PX \in \left\{\frac 12 (11 \pm \sqrt{61})\right\}\] so $PQ = \frac{\sqrt{61}}{2}$, i.e. $PQ^2 = \frac{61}{4}$ and the answer is $61+4 = \boxed{065}$.

 





\textbf{Solution 5}

Connect $AQ,QB$, since $\angle{AO_1P}=\angle{AOB}=\angle{BO_2P}$, so $\angle{AQP}=\frac{\angle{AO_1P}}{2}=\angle{BQP}=\frac{\angle{BO_2P}}{2}, \angle{AQB}=\angle{AOB}$ then, so $A,O,Q,B$ are concyclic

 We let $\angle{AO_1P}=\angle{AOB}=\angle{BO_2P}=2\alpha$, it is clear that $\angle{BQP}=\alpha, \angle{O_1AP}=90^{\circ}-\alpha$, which leads to the conclusion $OQ\bot XY$ which tells $Q$ is the midpoint of $XY$

 Then it is clear, $XP\cdot PY=15, XP=\frac{11-\sqrt{61}}{2}, PQ=\frac{11}{2}-\frac{11-\sqrt{61}}{2}=\frac{\sqrt{61}}{2}$ , the answer is $\boxed{065}$

 


 


 





\textbf{Solution 6(lazy)}

\begin{verbatim}
size(8cm);pair O, A, B, P, O1, O2, Q, X, Y;O=(0, 0);A=dir(140); B=dir(40);P=(3A+5B)/8;O1=extension((A+P)/2, (A+P)/2+(0, 1), A, O);O2=extension((B+P)/2, (B+P)/2+(0, 1), B, O);Q=intersectionpoints(circle(O1, length(A-O1)), circle(O2, length(B-O2)))[1];X=intersectionpoint(P -- (P+(P-Q)*100), circle(O, 1));Y=intersectionpoint(Q -- (Q+(Q-P)*100), circle(O, 1));draw(circle(O, 1));draw(circle(O1, length(A-O1)));draw(circle(O2, length(B-O2)));draw(A -- B,red); draw(X -- Y,green);dot("$A$", A, A);dot("$B$", B, B);dot("$P$", P, dir(70),blue);dot("$Q$", Q, dir(200));dot("$X$", X, 676767);dot("$Y$", Y, Y);label("$3$", (A+P)/2, N, red);label("$5$", (B+P)/2, N, red);draw(brace(X,Y));label("$11$",brace(X,Y),dir(20)); (Error making remote request. Unknown error\_msg)
\end{verbatim}$PX \cdot PY=AP \cdot PB=5 \cdot 3=15$ by power of a point. Also, $PX+PY=XY=11$, so $PX$ and $PY$ are solutions to the quadratic $x^2-11x+15=0$ so $PX$ and $PY$ is $\frac{11\pm\sqrt{61}}{2}$ in some order. Now, because we want $PQ^2$ and it is known to be rational, we can guess that $PQ$ is irrational or the problem would simply ask for $PQ$. We can also figure out that since $PQ^2$ is rational, $PQ$ is $\sqrt{\text{[something]}}$. $PQ=QX-PX$, and chances are low that $QX$ is some number with a square root plus or minus $\frac{\sqrt{61}}{2}$ to cancel out the $\frac{\sqrt{61}}{2}$ in $PX$, so one can see that $PQ^2$ is most likely to be $\left(\frac{\sqrt{61}}{2}\right)^2=\frac{61}{4}$, and our answer is $61+4=\boxed{065}$

 Note : If our answer is correct, then $QX=\frac{11}{2}$, which made $Q$ the midpoint of $XY$, a feature that occurs often in AIME problems, so that again made our answer probable. Midpoints have many properties and there is a lot of ways to show if a point is the midpoint of a segment. Even if the answer is wrong, it's still the same as leaving it blank and 065 is a good guess. ~\href{https://artofproblemsolving.com/wiki/index.php?title=Ddk001}{Ddk001}

 





\textbf{Solution 7}

We will show that $Q$ is the midpoint of $XY.$ To do this, let $Q^{\prime}$ be the altitude from $O$ to $XY$ or, equivalently, to $PQ.$ Notice that $O_{1}00_{2}P$ is a parallelogram. Thus, the height from $O$ to $O_{1}O_{2}$ is equal to the height from $P$ to $O_{1}O_{2}.$ Say that the line through $P$ perpendicular to $O_{1}O_{2}$ intersects $O_{1}O_{2}$ at $H.$ Then, $PQ$ is perpendicular to $O_{1}O_{2},$ so $H$ is on $PQ.$ Now, we have that the altitude from $O$ to $O_{1}O_{2}$ is equal to the altitude from $Q^{\prime}$ to $O_{1}O_{2}$ (since $OQ^{\prime} \parallel O_{1}O_{2}$). However, the altitude from $Q^{\prime}$ to $O_{1}O_{2}$ is just $Q^{\prime}H.$ Also, the altitude from $P$ to $O_{1}O_{2}$ is $PH$, so $PH = Q^{\prime}H.$ Thus, $O_{1}O_{2}$ bisects $PQ^{\prime}.$ However, this is true for $Q,$ too, so $Q = Q^{\prime},$ and we are done. Now, by PoP, we have 
 \[AP \cdot BP = XP \cdot YP = 15.\]Also, we have $XY = XP+YP = 11,$ so $XP = \frac{11 \pm \sqrt{61}}{2}$. Notice that $XQ = \frac{XY}{2} = \frac{11}{2},$ so $PQ = \frac{\sqrt{61}}{2},$ giving us our answer of $\boxed{065}.$

 


 





\textbf{Solution 8}


                    \begin{figure}[H]
                        \centering
                        \includegraphics[width=0.5\textwidth]{Solution/AIME_2019_I_15_Solution_Figure_3.png}
                    \end{figure}
                    \href{https://www.overleaf.com/read/bkdngbhwcskc\#415c3f}{Latex}Like Solution 7, let $Q'$ be the altitude from $O$ to $XY$. And, let $M$ be the intersection of $O_1O_2$ and $PQ$. Construct $P'$ on line $AO$ such that $PP' \parallel O_2O_1$. First, because of isosceles triangles $OAB$, $O_1AP$, and $O_2BP$, we have $\angle{OAP} = \angle{OBA} = \angle{APO_1} = \angle{BPO_2}$, which means $OO_1PO_2$ is a parallelogram. So, $O_2P = OO_1$. It is also clear that $PP'O_1O_2$ is a parallelogram by virtue of our definition. Thus, $O_2P = O_1P' = OO_1$. Since $OQ' \parallel O_1O_2 \parallel P'P$ (because of the right angles), $\frac{Q'M}{MP} = \frac{OO_1}{O_1P'} = 1 \implies Q'M = MP$. And, because $QM = MP$, $Q = Q'$. From Power of a Point on $P$, we have $XP(11-XP) = 15$, giving us $XP = \frac{11 - \sqrt{61}}{2}$. Since $OQ$ is perpendicular to $XY$, $Q$ is the midpoint of $XY$, so $XQ = \frac{11}{2}$. Thus, $PQ = \frac{11}{2} - \frac{11 - \sqrt{61}}{2} = \frac{\sqrt{61}}{2} \implies {PQ}^2 = \frac{61}{4}$. Therefore, our answer is $\boxed{65}$.

 $\newline$

 




\end{enumerate}